% ==============================================================================
% CHƯƠNG 2: TÌM HIỂU BÀI TOÁN VÀ PHÂN TÍCH THIẾT KẾ HỆ THỐNG
% ==============================================================================
\chapter{TÌM HIỂU BÀI TOÁN VÀ PHÂN TÍCH THIẾT KẾ HỆ THỐNG}
\label{chap:phan_tich_thiet_ke}

\section{Khảo sát nghiệp vụ và xác định yêu cầu}
\label{sec:khao_sat_nghiep_vu}

\subsection{Mục đích và mục tiêu của hệ thống}
Mục đích cốt lõi của đề tài là xây dựng một hệ thống phần mềm quản lý bán lẻ tại điểm bán (Point of Sale - POS) dựa trên nền tảng Web hiện đại, giải quyết toàn diện các bài toán thường gặp của chuỗi cửa hàng bán lẻ và dịch vụ ẩm thực (F\&B):
\begin{enumerate}[label=\arabic*., leftmargin=1.5cm]
    \item \textbf{Tối ưu hóa tốc độ phục vụ tại quầy:} Rút ngắn thời gian lên đơn, tìm kiếm sản phẩm và thanh toán xuống dưới $1$ giây cho mỗi thao tác, giảm thiểu tối đa hiện tượng ùn tắc khách hàng trong các khung giờ cao điểm.
    \item \textbf{Tự động hóa thanh toán không dùng tiền mặt:} Tích hợp thanh toán quét mã VietQR động theo thời gian thực kết nối với cổng PayOS, loại bỏ hoàn toàn các sai sót khi nhập thủ công số tiền hoặc số tài khoản.
    \item \textbf{Quản lý kho hàng theo nguyên tắc kế toán kép (Ledger Inventory):} Ghi nhận toàn bộ biến động nhập, xuất, hao hụt, kiểm kê vào sổ cái giao dịch, duy trì lịch sử biến động minh bạch và ngăn ngừa gian lận thất thoát.
    \item \textbf{Hỗ trợ các chương trình khuyến mãi linh hoạt:} Tự động đánh giá và áp dụng mức giảm giá tối ưu nhất cho khách hàng dựa trên các điều kiện về khung giờ vàng (Happy Hour), mua hàng tặng hàng (BOGO) hoặc mã giảm giá (Coupon), tuân thủ nguyên tắc không cộng dồn khuyến mãi (No Stacking).
\end{enumerate}

\subsection{Khảo sát hiện trạng và quy trình bán hàng tại quầy}
Khảo sát thực tế tại các điểm bán lẻ truyền thống cho thấy quy trình bán hàng thủ công tồn tại nhiều điểm nghẽn nghiêm trọng:
\begin{itemize}[leftmargin=1.5cm]
    \item Khi khách hàng chọn món có nhiều biến thể (ví dụ: Áo thun cùng mẫu nhưng khác Size M/L/XL và khác Màu sắc; hoặc Trà sữa có mức đường, đá và topping khác nhau), nhân viên thường phải ghi chú viết tay vào phiếu giấy, dẫn đến sai sót khi chế biến hoặc xuất kho.
    \item Thu ngân phải tự tính nhẩm tiền thừa hoặc áp dụng mã giảm giá thủ công, rất dễ gây nhầm lẫn hoặc thâm hụt ngân quỹ vào cuối ca làm việc.
    \item Việc kiểm tra số lượng tồn kho chỉ được thực hiện vào cuối ngày hoặc cuối tháng bằng cách đếm thủ công, không có khả năng cảnh báo kịp thời khi một mặt hàng bán chạy sắp hết hàng trong ca bán.
\end{itemize}

Hệ thống POS xây dựng trong đề tài số hóa toàn diện quy trình này: Toàn bộ danh mục sản phẩm, biến thể (Variant SKU) và nhóm tùy chọn bổ sung (Modifier Groups) được đồng bộ hóa lên màn hình cảm ứng; việc tính toán tổng tiền, thuế VAT (10\%) và khuyến mãi được máy chủ thực hiện tự động và trả kết quả tức thì.

\subsection{Phạm vi hệ thống và đối tượng sử dụng}
Hệ thống được thiết kế hướng tới phục vụ 2 nhóm đối tượng người dùng chính (Actors):
\begin{enumerate}[label=\arabic*., leftmargin=1.5cm]
    \item \textbf{Nhân viên thu ngân (Cashier / Role Staff):}
    \begin{itemize}
        \item Thực hiện bán hàng nhanh tại quầy thông qua màn hình \textit{Explore POS}.
        \item Tìm kiếm sản phẩm theo tên, danh mục hoặc mã vạch SKU.
        \item Tùy biến biến thể (Size, Màu sắc) và các nhóm tùy chọn bổ sung (Topping, Đá, Đường).
        \item Tra cứu nhanh hoặc đăng ký mới khách hàng thân thiết bằng số điện thoại.
        \item Nhập mã khuyến mãi hoặc để hệ thống tự động áp dụng chương trình ưu đãi tốt nhất.
        \item Thực hiện thanh toán Tiền mặt (tự động tính tiền thừa và in hóa đơn nhiệt 80mm) hoặc chuyển khoản VietQR PayOS.
        \item Chuyển đổi phương thức thanh toán sang tiền mặt khi chuyển khoản gặp sự cố; tra cứu lịch sử hóa đơn trong ca và kích hoạt hủy đơn hàng chờ (kèm hoàn kho tự động).
    \end{itemize}
    \item \textbf{Quản lý / Quản trị viên (Store Manager / Admin):}
    \begin{itemize}
        \item Quản lý toàn bộ danh mục sản phẩm, biến thể phân loại động (JSON attributes) và nhóm tùy chọn.
        \item Nhập kho, xuất kho hao hụt và thực hiện kiểm kê cân bằng tồn kho thực tế.
        \item Thiết lập các chương trình khuyến mãi đa dạng (COUPON, HAPPY\_HOUR, BOGO).
        \item Quản lý danh bạ khách hàng CRM, theo dõi tích lũy chi tiêu và số lượng đơn hàng.
        \item Quản lý tài khoản nhân viên, cấp phát vai trò và quyền hạn bảo mật RBAC.
        \item Theo dõi bảng điều khiển chỉ số kinh doanh (Dashboard) và kiểm toán dấu vết thao tác qua Activity Logs.
    \end{itemize}
\end{enumerate}

\textit{\textbf{Ghi chú về vai trò Thủ kho:} Trong phạm vi đề tài, chức năng quản lý kho hàng (nhập kho, xuất kho, kiểm kê) được tích hợp trực tiếp vào vai trò Quản trị viên do quy mô cửa hàng nhỏ và vừa không yêu cầu phân tách tổ chức kho riêng biệt. Quản trị viên đồng thời đảm nhiệm nghiệp vụ của Thủ kho, giúp đơn giản hóa phân quyền RBAC và phù hợp với mô hình vận hành thực tế.}

\subsection{Yêu cầu chức năng chi tiết (Functional Requirements)}
Hệ thống được phân rã thành 6 phân hệ chức năng chính, như trình bày tại Bảng \ref{tab:yeu_cau_chuc_nang}.

\begin{table}[htbp]
\centering
\caption{Bảng đặc tả các yêu cầu chức năng hệ thống POS}
\label{tab:yeu_cau_chuc_nang}
\vspace{2mm}
\begin{tabularx}{\textwidth}{|c|p{3.5cm}|X|c|}
\hline
\textbf{STT} & \textbf{Phân hệ} & \textbf{Mô tả chi tiết chức năng} & \textbf{Tác nhân} \\ \hline
1 & \textbf{Xác thực \& Bảo mật (Auth)} & Đăng nhập tài khoản bằng email/mật khẩu, cấp Access Token ngắn hạn và HttpOnly Refresh Token dài hạn, tự động xoay vòng token khi hết hạn, đăng xuất và thu hồi phiên. & Tất cả \\ \hline
2 & \textbf{Bán hàng POS (Checkout)} & Tìm kiếm sản phẩm, lọc theo danh mục, chọn biến thể/modifier, quản lý giỏ hàng, nhận diện khách hàng CRM, tính toán chiết khấu/VAT, in hóa đơn nhiệt. & Thu ngân \\ \hline
3 & \textbf{Thanh toán đa kênh (Payment)} & Thanh toán tiền mặt (tính tiền thừa), thanh toán trực tuyến qua VietQR PayOS động, tự động nhận Webhook cập nhật đơn, cho phép chuyển sang tiền mặt tại quầy. & Thu ngân, PayOS \\ \hline
4 & \textbf{Quản lý Kho sổ cái (Inventory)} & Quản lý tồn kho dựa trên biến thể SKU, ghi nhận giao dịch sổ cái (IN/OUT/ADJUSTMENT), kiểm kê kho thực tế và tự động bù trừ kho khi hủy đơn. & Quản lý, Thu ngân \\ \hline
5 & \textbf{Khuyến mãi (Promotions)} & Thiết lập 3 loại ưu đãi: COUPON (nhập mã), HAPPY\_HOUR (khung giờ vàng theo ngày/thứ), BOGO (mua A tặng B). Đánh giá tự động no-stacking. & Quản lý, Thu ngân \\ \hline
6 & \textbf{Quản trị \& Kiểm toán (Admin \& Logs)} & Quản lý sản phẩm/danh mục/người dùng/khách hàng, theo dõi Dashboard doanh thu và tra cứu nhật ký kiểm toán hoạt động Activity Logs. & Quản trị viên \\ \hline
\end{tabularx}
\end{table}

\subsection{Yêu cầu phi chức năng (Non-Functional Requirements)}
Để đảm bảo hệ thống vận hành ổn định trong môi trường bán lẻ thực tế với cường độ cao, các yêu cầu phi chức năng được xác định nghiêm ngặt:
\begin{enumerate}[label=\arabic*., leftmargin=1.5cm]
    \item \textbf{Hiệu năng và Thời gian phản hồi (Performance):}
    \begin{itemize}
        \item Thời gian thực thi thao tác thêm sản phẩm vào giỏ hàng hoặc tính toán lại khuyến mãi trên giao diện Front-end phải nhỏ hơn $50$ mili-giây.
        \item Thời gian đáp ứng của máy chủ API đối với các yêu cầu bán hàng, tra cứu thông tin sản phẩm và khách hàng phải đạt dưới $1$ giây trong điều kiện mạng tiêu chuẩn.
    \end{itemize}
    \item \textbf{Tính toàn vẹn tài chính và Kế toán (Financial Integrity - ADR-0006):}
    \begin{itemize}
        \item Nghiêm cấm hoàn toàn hành vi xóa các đơn hàng đã ở trạng thái hoàn tất (\texttt{COMPLETED}) nhằm bảo toàn tính toàn vẹn của sổ sách kế toán và báo cáo thuế doanh nghiệp.
        \item Khi một đơn hàng đang ở trạng thái chờ (\texttt{PENDING}) bị hủy bỏ hoặc xóa, hệ thống bắt buộc phải tự động kích hoạt quy trình giao dịch bù trừ toàn diện (Compensating Ledger Transaction): hoàn lại số lượng hàng hóa vào kho (giao dịch \texttt{IN}), hoàn lại số lượt sử dụng voucher khuyến mãi, hoàn lại số liệu doanh số CRM khách hàng và gửi lệnh hủy liên kết thanh toán sang cổng PayOS.
    \end{itemize}
    \item \textbf{Tính bảo mật và An toàn thông tin (Security - ADR-0007):}
    \begin{itemize}
        \item Mật khẩu người dùng bắt buộc phải được băm bảo mật bằng thuật toán BCrypt với hệ số Salt tối thiểu là 10.
        \item Tuyệt đối không lưu trữ khóa xác thực JWT trong bộ nhớ cục bộ \texttt{localStorage} của trình duyệt; áp dụng cơ chế Cookie HttpOnly kết hợp phát hiện tấn công tái sử dụng Refresh Token (Replay Attack Detection).
        \item Mọi thao tác làm thay đổi dữ liệu nhạy cảm (CREATE, UPDATE, DELETE) đều phải được ghi nhận tự động vào bảng nhật ký kiểm toán \texttt{tbl\_activity\_logs}.
    \end{itemize}
    \item \textbf{Tính khả dụng và Khả năng phục vụ liên tục (Availability):} Giao diện hệ thống tương thích tốt trên các kích thước màn hình phổ biến từ máy tính bảng 10 inch đến màn hình máy tính Full HD; hỗ trợ chế độ in nhiệt chuẩn 80mm/58mm.
\end{enumerate}

% ==============================================================================
\section{Phân tích và thiết kế chức năng}
\label{sec:thiet_ke_chuc_nang}

\subsection{Biểu đồ phân rã chức năng (BFD - Business Function Diagram)}
Biểu đồ phân rã chức năng (Hình \ref{fig:bfd_diagram}) mô tả cấu trúc phân cấp toàn bộ cây nghiệp vụ của hệ thống POS từ cấp hệ thống đến các mô-đun chức năng chi tiết.

\begin{figure}[H]
\centering
\begin{minipage}{0.95\textwidth}
\centering
\fbox{
    \parbox{0.95\textwidth}{
        \centering \vspace{2mm}
        {\fontsize{13pt}{15pt}\selectfont \textbf{HỆ THỐNG POS QUẢN LÝ BÁN HÀNG (BILLING SOFTWARE)}}\\[3mm]
        \rule{0.9\textwidth}{0.8pt}\\[2mm]
        \begin{tabularx}{\textwidth}{XXXX}
            \textbf{1. BÁN HÀNG TẠI QUẦY} & \textbf{2. QUẢN LÝ KHO HÀNG} & \textbf{3. QUẢN LÝ DANH MỤC} & \textbf{4. QUẢN TRỊ HỆ THỐNG} \\
            $\bullet$ Tìm kiếm hàng hóa & $\bullet$ Nhập kho (Stock IN) & $\bullet$ Quản lý Danh mục & $\bullet$ Quản lý Người dùng \\
            $\bullet$ Chọn Variant/Modifier & $\bullet$ Xuất kho (Stock OUT) & $\bullet$ Sản phẩm \& Biến thể & $\bullet$ Bảng điều khiển KPI \\
            $\bullet$ Áp dụng Khuyến mãi & $\bullet$ Kiểm kê cân bằng kho & $\bullet$ Nhóm tùy chọn Modifier & $\bullet$ Lịch sử kiểm toán Log \\
            $\bullet$ Thanh toán tiền mặt & $\bullet$ Lịch sử sổ cái kho & $\bullet$ Thiết lập Khuyến mãi & $\bullet$ Quản lý phiên token \\
            $\bullet$ Thanh toán VietQR & $\bullet$ Cảnh báo hết hàng & $\bullet$ Danh bạ khách hàng CRM & $\bullet$ Đồng bộ Webhook \\
            $\bullet$ In hóa đơn nhiệt & $\bullet$ Bù trừ hủy đơn & $\bullet$ Tải ảnh đám mây S3 & $\bullet$ Phân quyền RBAC \\
        \end{tabularx}\\[2mm]
    }
}
\end{minipage}
\caption{Biểu đồ phân rã chức năng (BFD) của hệ thống POS}
\label{fig:bfd_diagram}
\end{figure}

\subsection{Mô hình Use Case}

\subsubsection{Biểu đồ Use Case tổng quát}
Hệ thống bao gồm 2 tác nhân con người (\textit{Thu ngân} và \textit{Quản lý}) cùng 1 tác nhân hệ thống bên thứ ba (\textit{Cổng thanh toán PayOS}). Mối quan hệ giữa các tác nhân và các ca sử dụng tổng quát được minh họa tại Hình \ref{fig:usecase_overall}.

\begin{figure}[H]
\centering
\begin{minipage}{0.95\textwidth}
\centering
\fbox{
    \parbox{0.95\textwidth}{
        \centering \vspace{2mm}
        \textbf{BIỂU ĐỒ USE CASE TỔNG QUÁT TOÀN HỆ THỐNG}\\[3mm]
        \begin{tabularx}{\textwidth}{p{3.5cm}Xp{3.5cm}}
            \centering \textbf{ACTOR: THU NGÂN} & \centering \textbf{CÁC USE CASE HỆ THỐNG} & \centering \textbf{ACTOR: QUẢN LÝ} \tabularnewline
            $\bullet$ Bán hàng POS tại quầy & $\Longleftrightarrow$ [UC01: Đăng nhập \& Phiên làm việc] $\Longleftrightarrow$ & $\bullet$ Toàn quyền quản trị \tabularnewline
            $\bullet$ Tra cứu thông tin món & $\Longleftrightarrow$ [UC02: Tạo đơn \& Bán lẻ POS] & $\bullet$ Quản lý sản phẩm \tabularnewline
            $\bullet$ Chọn biến thể / Topping & $\Longleftrightarrow$ [UC03: Đánh giá Khuyến mãi] $\Longleftrightarrow$ & $\bullet$ Quản lý khuyến mãi \tabularnewline
            $\bullet$ Nhận diện khách hàng & $\Longleftrightarrow$ [UC04: Thanh toán VietQR / Cash] & $\bullet$ Nhập / Kiểm kê kho \tabularnewline
            $\bullet$ Chuyển sang tiền mặt & $\Longleftrightarrow$ [UC05: Quản lý Kho sổ cái] $\Longleftrightarrow$ & $\bullet$ Xem Dashboard KPI \tabularnewline
            $\bullet$ In lại hóa đơn bán lẻ & $\Longleftrightarrow$ [UC06: Tra cứu Activity Logs] $\Longleftrightarrow$ & $\bullet$ Quản lý nhân viên \tabularnewline
            & $\Longleftrightarrow$ [UC07: Xem Báo cáo Doanh thu] $\Longleftrightarrow$ & $\bullet$ Xem báo cáo thống kê \tabularnewline
        \end{tabularx}\\[2mm]
        \rule{0.9\textwidth}{0.5pt}\\[1mm]
        \textit{Actor bên ngoài: [PayOS Gateway] $\Longleftrightarrow$ Gửi Webhook xác thực đơn hàng thành công}\\[2mm]
    }
}
\end{minipage}
\caption{Biểu đồ Use Case tổng quát hệ thống POS}
\label{fig:usecase_overall}
\end{figure}

\subsubsection{Đặc tả chi tiết 7 Use Cases cốt lõi}

Dưới đây là các bảng đặc tả chuẩn học thuật cho 7 ca sử dụng cốt lõi của hệ thống POS:

% ------------------------------------------------------------------------------
% USE CASE 01
% ------------------------------------------------------------------------------
\begin{table}[htbp]
\centering
\caption{Bảng đặc tả chi tiết Use Case UC01: Đăng nhập và Khởi tạo phiên làm việc}
\label{tab:uc01_login}
\vspace{2mm}
\begin{tabularx}{\textwidth}{|l|X|}
\hline
\textbf{Mã Use Case} & \textbf{UC01} \\ \hline
\textbf{Tên Use Case} & Đăng nhập và Khởi tạo phiên làm việc bảo mật \\ \hline
\textbf{Tác nhân (Actor)} & Nhân viên thu ngân, Quản trị viên \\ \hline
\textbf{Mô tả tóm tắt} & Người dùng nhập thông tin xác thực để truy cập vào hệ thống, nhận Access Token trong bộ nhớ và Cookie Refresh Token an toàn. \\ \hline
\textbf{Tiền điều kiện} & Tài khoản người dùng đã được quản trị viên cấp phát trên hệ thống và đang ở trạng thái hoạt động. \\ \hline
\textbf{Luồng sự kiện chính} & 
\begin{enumerate}[leftmargin=*, nosep]
    \item Người dùng truy cập đường dẫn \texttt{/login} trên trình duyệt.
    \item Hệ thống hiển thị biểu mẫu yêu cầu nhập Email và Mật khẩu.
    \item Người dùng nhập thông tin và nhấn nút ''Đăng nhập''.
    \item Front-end gửi yêu cầu \texttt{POST /api/v1.0/login} kèm thông tin.
    \item Back-end kiểm tra tài khoản, đối chiếu mật khẩu đã băm bằng BCrypt.
    \item Khi thông tin hợp lệ, Back-end sinh Access Token (hạn 15 phút) và Refresh Token ngẫu nhiên (hạn 7 ngày).
    \item Back-end băm SHA-256 Refresh Token và lưu vào bảng \texttt{refresh\_tokens}, ghi log đăng nhập vào \texttt{tbl\_activity\_logs}.
    \item Back-end thiết lập Cookie \texttt{refreshToken} với cờ \texttt{HttpOnly} và trả về DTO chứa Access Token, vai trò người dùng.
    \item Front-end lưu Access Token vào biến bộ nhớ RAM (\texttt{authSession.js}) và tự động điều hướng người dùng vào trang \texttt{/dashboard} hoặc \texttt{/explore}.
\end{enumerate} \\ \hline
\textbf{Luồng ngoại lệ} & 
\begin{itemize}[leftmargin=*, nosep]
    \item \textit{3a. Email không tồn tại hoặc mật khẩu sai:} Hệ thống trả mã lỗi 401 Unauthorized, Front-end hiển thị thông báo lỗi bằng tiếng Việt: ''Email hoặc mật khẩu không chính xác''.
    \item \textit{3b. Lỗi kết nối máy chủ:} Front-end thông báo: ''Không thể kết nối đến máy chủ API''.
\end{itemize} \\ \hline
\textbf{Hậu điều kiện} & Người dùng có phiên làm việc hợp lệ; mọi request tiếp theo đều được đính kèm Access Token trong tiêu đề Authorization. \\ \hline
\end{tabularx}
\end{table}

% ------------------------------------------------------------------------------
% USE CASE 02
% ------------------------------------------------------------------------------
\begin{table}[htbp]
\centering
\caption{Bảng đặc tả chi tiết Use Case UC02: Tạo đơn hàng và Bán lẻ tại quầy POS}
\label{tab:uc02_pos}
\vspace{2mm}
\begin{tabularx}{\textwidth}{|l|X|}
\hline
\textbf{Mã Use Case} & \textbf{UC02} \\ \hline
\textbf{Tên Use Case} & Tạo đơn hàng và Bán lẻ tại quầy POS \\ \hline
\textbf{Tác nhân (Actor)} & Nhân viên thu ngân \\ \hline
\textbf{Mô tả tóm tắt} & Thu ngân chọn danh mục, tùy biến biến thể hàng hóa, quản lý giỏ hàng, nhận diện khách hàng và xác nhận tạo đơn hàng. \\ \hline
\textbf{Tiền điều kiện} & Thu ngân đã đăng nhập thành công vào hệ thống và đang mở màn hình \texttt{/explore}. \\ \hline
\textbf{Luồng sự kiện chính} & 
\begin{enumerate}[leftmargin=*, nosep]
    \item Thu ngân lọc danh mục hoặc gõ từ khóa tìm kiếm mặt hàng.
    \item Thu ngân nhấp vào sản phẩm:
    \begin{itemize}
        \item Nếu sản phẩm có nhiều biến thể hoặc tùy chọn bổ sung: Hệ thống mở popup \texttt{POSItemModal}, thu ngân chọn kích thước/màu sắc (Variant) và đánh dấu các tùy chọn Topping/Đá/Đường (Modifiers).
        \item Nếu sản phẩm đơn giản: Hệ thống thêm trực tiếp vào giỏ hàng.
    \end{itemize}
    \item Thu ngân điều chỉnh số lượng hoặc xóa mặt hàng trong giỏ \texttt{CartItems}.
    \item Thu ngân nhập số điện thoại khách hàng vào \texttt{CustomerForm}; hệ thống tự động tìm kiếm và hiển thị tên, điểm tích lũy của khách hàng thân thiết (nếu khách hàng mới, thu ngân có thể tạo nhanh).
    \item Thu ngân nhấn nút ''Thanh toán'', chọn hình thức thanh toán (Tiền mặt hoặc VietQR).
    \item Front-end gửi yêu cầu \texttt{POST /api/v1.0/orders} chứa danh sách sản phẩm, biến thể, modifiers và thông tin khách hàng.
    \item Back-end tiếp nhận, gọi \texttt{OrderService} tính toán lại giá tiền phía server, ghi giao dịch xuất kho \texttt{OUT} vào bảng sổ cái tồn kho, cập nhật số dư kho, và lưu đơn hàng vào bảng \texttt{tbl\_orders}.
\end{enumerate} \\ \hline
\textbf{Luồng ngoại lệ} & 
\begin{itemize}[leftmargin=*, nosep]
    \item \textit{2a. Mặt hàng đã hết hàng trong kho:} Hệ thống hiển thị cảnh báo tồn kho không đủ (nếu không bật chế độ bán âm).
    \item \textit{4a. Số điện thoại khách hàng không đúng định dạng:} Front-end hiển thị lỗi kiểm tra dữ liệu yêu cầu nhập đủ 10 chữ số.
\end{itemize} \\ \hline
\textbf{Hậu điều kiện} & Đơn hàng mới được ghi nhận vào cơ sở dữ liệu; số lượng tồn kho của các biến thể liên quan bị trừ tương ứng trong sổ cái kế toán. \\ \hline
\end{tabularx}
\end{table}

% ------------------------------------------------------------------------------
% USE CASE 03
% ------------------------------------------------------------------------------
\begin{table}[htbp]
\centering
\caption{Bảng đặc tả chi tiết Use Case UC03: Đánh giá và Áp dụng Khuyến mãi}
\label{tab:uc03_promotion}
\vspace{2mm}
\begin{tabularx}{\textwidth}{|l|X|}
\hline
\textbf{Mã Use Case} & \textbf{UC03} \\ \hline
\textbf{Tên Use Case} & Đánh giá và Áp dụng chương trình Khuyến mãi \\ \hline
\textbf{Tác nhân (Actor)} & Thu ngân, Hệ thống tính toán khuyến mãi \\ \hline
\textbf{Mô tả tóm tắt} & Hệ thống tự động phân tích giỏ hàng, đối chiếu các chương trình ưu đãi hiện hành và chọn ra 1 khuyến mãi có mức giảm giá cao nhất (No Stacking). \\ \hline
\textbf{Tiền điều kiện} & Giỏ hàng có ít nhất một mặt hàng hợp lệ. \\ \hline
\textbf{Luồng sự kiện chính} & 
\begin{enumerate}[leftmargin=*, nosep]
    \item Mỗi khi giỏ hàng thay đổi (thêm/bớt mặt hàng, nhập mã giảm giá Coupon): Front-end gọi API \texttt{POST /api/v1.0/promotions/evaluate}.
    \item Back-end lấy danh sách tất cả khuyến mãi đang ở trạng thái kích hoạt (\texttt{isActive = true}) trong cơ sở dữ liệu.
    \item Bộ máy đánh giá (\texttt{PromotionServiceImpl}) tiến hành kiểm tra:
    \begin{itemize}
        \item Khuyến mãi \texttt{HAPPY\_HOUR}: Kiểm tra ngày trong tuần (\texttt{daysOfWeek}), khung thời gian bắt đầu/kết thúc (\texttt{startTime}/\texttt{endTime}) và giá trị đơn hàng tối thiểu (\texttt{minOrderAmount}).
        \item Khuyến mãi \texttt{BOGO}: Kiểm tra số lượng cặp biến thể mua (\texttt{buyVariantId}) và tặng (\texttt{getVariantId}) trong giỏ hàng.
        \item Khuyến mãi \texttt{COUPON}: So sánh mã người dùng nhập, kiểm tra thời hạn và giới hạn số lượt sử dụng (\texttt{usageLimit} so với \texttt{timesUsed}).
    \end{itemize}
    \item Hệ thống tính số tiền giảm giá của từng chương trình hợp lệ, chọn ra duy nhất 1 chương trình có số tiền giảm cao nhất.
    \item Back-end trả về kết quả DTO gồm: mã khuyến mãi được chọn, tên chương trình, số tiền chiết khấu thực tế (\texttt{discountAmount}) và tổng tiền mới.
    \item Front-end cập nhật số tiền giảm giá và thanh toán lên giao diện giỏ hàng.
\end{enumerate} \\ \hline
\textbf{Luồng ngoại lệ} & 
\begin{itemize}[leftmargin=*, nosep]
    \item \textit{3a. Mã Coupon không tồn tại hoặc đã hết hạn/hết lượt dùng:} Back-end trả về thông báo lỗi cụ thể, giỏ hàng giữ nguyên giá trị ban đầu và không áp dụng chiết khấu.
\end{itemize} \\ \hline
\textbf{Hậu điều kiện} & Hóa đơn được áp dụng đúng mức chiết khấu tối ưu nhất; thông tin khuyến mãi được gắn kèm vào đơn hàng khi tạo mới. \\ \hline
\end{tabularx}
\end{table}

% ------------------------------------------------------------------------------
% USE CASE 04
% ------------------------------------------------------------------------------
\begin{table}[htbp]
\centering
\caption{Bảng đặc tả chi tiết Use Case UC04: Thanh toán VietQR PayOS và Chuyển đổi tại quầy}
\label{tab:uc04_payos}
\vspace{2mm}
\begin{tabularx}{\textwidth}{|l|X|}
\hline
\textbf{Mã Use Case} & \textbf{UC04} \\ \hline
\textbf{Tên Use Case} & Thanh toán trực tuyến VietQR PayOS và Chuyển đổi phương thức tại quầy \\ \hline
\textbf{Tác nhân (Actor)} & Thu ngân, Khách hàng, Cổng thanh toán PayOS \\ \hline
\textbf{Mô tả tóm tắt} & Sinh mã QR động cho khách hàng quét thanh toán ngân hàng, tự động nhận Webhook hoàn tất đơn; hỗ trợ chuyển sang tiền mặt khi gặp sự cố mạng. \\ \hline
\textbf{Tiền điều kiện} & Đơn hàng đã được tạo thành công với phương thức thanh toán là \texttt{PAYOS} và trạng thái là \texttt{PENDING}. \\ \hline
\textbf{Luồng sự kiện chính} & 
\begin{enumerate}[leftmargin=*, nosep]
    \item Hệ thống hiển thị cửa sổ \texttt{PaymentQRCode} chứa mã VietQR động được tạo từ PayOS SDK, hiển thị chính xác tổng tiền và mã đơn.
    \item Front-end kích hoạt cơ chế thăm dò (Polling) qua hook \texttt{usePolledOrderData} với chu kỳ $3$ giây/lần.
    \item Khách hàng mở ứng dụng ngân hàng di động bất kỳ, quét mã VietQR và xác nhận chuyển tiền.
    \item Ngân hàng ghi có vào tài khoản cửa hàng; hệ thống PayOS gửi ngay một tín hiệu Webhook (\texttt{POST /api/v1.0/payos/webhook}) tới máy chủ Back-end.
    \item Back-end xác thực tính hợp lệ của chữ ký số HMAC-SHA256 Checksum.
    \item Back-end cập nhật trạng thái đơn hàng sang \texttt{COMPLETED}, cập nhật điểm tích lũy CRM của khách hàng.
    \item Tại lần thăm dò tiếp theo, Front-end nhận trạng thái \texttt{COMPLETED}, tự động đóng popup QR và mở mẫu in hóa đơn bán lẻ \texttt{ReceiptPopup}.
\end{enumerate} \\ \hline
\textbf{Luồng ngoại lệ (Quy trình chuyển đổi thanh toán - ADR-0006):} & 
\begin{itemize}[leftmargin=*, nosep]
    \item \textit{3a. Khách hàng không thể chuyển khoản (mất mạng, lỗi app):} 
    \begin{enumerate}[label=\alph*., leftmargin=*]
        \item Thu ngân nhấn nút ''Chuyển sang tiền mặt'' trên màn hình thanh toán.
        \item Front-end gửi yêu cầu \texttt{POST /api/v1.0/orders/\{orderId\}/switch-to-cash}.
        \item Back-end đổi phương thức thanh toán sang \texttt{CASH}, cập nhật trạng thái đơn thành \texttt{COMPLETED} và gọi API PayOS để hủy link thanh toán từ xa.
        \item Hệ thống hoàn tất giao dịch mà không tạo đơn trùng lặp và không nhân đôi việc trừ số lượng kho.
    \end{enumerate}
\end{itemize} \\ \hline
\textbf{Hậu điều kiện} & Đơn hàng được thanh toán thành công, liên kết VietQR đóng lại và hóa đơn bán lẻ sẵn sàng để in. \\ \hline
\end{tabularx}
\end{table}

% ------------------------------------------------------------------------------
% USE CASE 05
% ------------------------------------------------------------------------------
\begin{table}[htbp]
\centering
\caption{Bảng đặc tả chi tiết Use Case UC05: Quản lý Biến thể và Sổ cái Kho}
\label{tab:uc05_inventory}
\vspace{2mm}
\begin{tabularx}{\textwidth}{|l|X|}
\hline
\textbf{Mã Use Case} & \textbf{UC05} \\ \hline
\textbf{Tên Use Case} & Quản lý Biến thể hàng hóa và Kiểm kê kho theo sổ cái \\ \hline
\textbf{Tác nhân (Actor)} & Quản trị viên, Quản lý kho \\ \hline
\textbf{Mô tả tóm tắt} & Quản lý thông tin các biến thể sản phẩm theo SKU, thực hiện nhập kho, xuất hủy và kiểm kê cân bằng số dư tồn kho thực tế. \\ \hline
\textbf{Tiền điều kiện} & Người dùng có quyền \texttt{ROLE\_ADMIN} đã đăng nhập vào hệ thống. \\ \hline
\textbf{Luồng sự kiện chính (Quy trình kiểm kê kho thực tế):} & 
\begin{enumerate}[leftmargin=*, nosep]
    \item Quản lý mở trang \texttt{/items}, chọn một sản phẩm và mở modal \texttt{StockOperationModal} của một biến thể (Variant SKU).
    \item Hệ thống hiển thị số lượng tồn kho sổ cái hiện tại (\texttt{cachedStockQuantity}) và lịch sử các giao dịch gần nhất.
    \item Quản lý chọn tab ''Kiểm kê kho'' (Stock Check) và nhập số lượng thực tế kiểm đếm được tại quầy.
    \item Hệ thống tự động tính độ lệch tồn kho: $\Delta = \text{Số thực tế} - \text{Số hệ thống}$.
    \item Quản lý nhập ghi chú kiểm kê và nhấn ''Xác nhận cân bằng kho''.
    \item Front-end gửi yêu cầu \texttt{POST /api/v1.0/admin/inventory/stock-check}.
    \item Back-end tạo một bản ghi giao dịch mới có kiểu \texttt{ADJUSTMENT} vào bảng \texttt{tbl\_inventory\_transactions} với số lượng là $\Delta$.
    \item Back-end tính lại tổng tồn kho từ toàn bộ giao dịch sổ cái và cập nhật trường \texttt{cachedStockQuantity} trong bảng \texttt{tbl\_variants}.
    \item Hệ thống ghi log kiểm toán vào \texttt{tbl\_activity\_logs}.
\end{enumerate} \\ \hline
\textbf{Luồng ngoại lệ} & 
\begin{itemize}[leftmargin=*, nosep]
    \item \textit{3a. Số lượng thực tế nhập vào là số âm:} Hệ thống hiển thị lỗi yêu cầu số lượng thực đếm phải lớn hơn hoặc bằng 0.
\end{itemize} \\ \hline
\textbf{Hậu điều kiện} & Số dư tồn kho trên hệ thống khớp chính xác với số lượng vật lý thực tế tại cửa hàng; toàn bộ sai lệch được lưu vết minh bạch trong sổ cái. \\ \hline
\end{tabularx}
\end{table}

% ------------------------------------------------------------------------------
% USE CASE 06
% ------------------------------------------------------------------------------
\begin{table}[htbp]
\centering
\caption{Bảng đặc tả chi tiết Use Case UC06: Tra cứu Lịch sử Kiểm toán Hoạt động}
\label{tab:uc06_activity_log}
\vspace{2mm}
\begin{tabularx}{\textwidth}{|l|X|}
\hline
\textbf{Mã Use Case} & \textbf{UC06} \\ \hline
\textbf{Tên Use Case} & Tra cứu Lịch sử Kiểm toán Hoạt động (Activity Logs) \\ \hline
\textbf{Tác nhân (Actor)} & Quản trị viên, Nhân viên thu ngân \\ \hline
\textbf{Mô tả tóm tắt} & Tra cứu và lọc lịch sử các thao tác dữ liệu nhạy cảm trong hệ thống nhằm mục đích đối soát an ninh và kiểm toán nội bộ. \\ \hline
\textbf{Tiền điều kiện} & Người dùng đã đăng nhập vào hệ thống. \\ \hline
\textbf{Luồng sự kiện chính} & 
\begin{enumerate}[leftmargin=*, nosep]
    \item Người dùng truy cập đường dẫn \texttt{/activity-logs}.
    \item Front-end gửi yêu cầu \texttt{GET /api/v1.0/activity-logs} kèm các tham số phân trang, bộ lọc hành động (\texttt{action}), loại thực thể (\texttt{entityType}) và khoảng thời gian.
    \item Back-end kiểm tra vai trò an ninh trong \texttt{SecurityContextHolder}:
    \begin{itemize}
        \item Nếu người dùng là \texttt{ROLE\_ADMIN}: Hệ thống trả về toàn bộ nhật ký thao tác của tất cả các tài khoản trong hệ thống.
        \item Nếu người dùng là \texttt{ROLE\_STAFF} (Thu ngân): Hệ thống tự động giới hạn chỉ trả về các bản ghi nhật ký do chính email của nhân viên đó thực hiện.
    \end{itemize}
    \item Giao diện hiển thị bảng nhật ký gồm: Thời gian thực hiện, Email người thực hiện, Loại hành động (LOGIN, CREATE, UPDATE, DELETE, CANCEL), Thực thể chịu tác động và Mô tả chi tiết.
\end{enumerate} \\ \hline
\textbf{Luồng ngoại lệ} & 
\begin{itemize}[leftmargin=*, nosep]
    \item \textit{2a. Khoảng thời gian lọc không hợp lệ (Ngày bắt đầu lớn hơn ngày kết thúc):} Hệ thống hiển thị cảnh báo lỗi dữ liệu nhập.
\end{itemize} \\ \hline
\textbf{Hậu điều kiện} & Người dùng nắm bắt được lịch sử biến động dữ liệu; mọi hành vi đáng ngờ đều có thể truy xuất chính xác người thực hiện và thời điểm. \\ \hline
\end{tabularx}
\end{table}

% ------------------------------------------------------------------------------
% USE CASE 07 - Xem báo cáo thống kê doanh thu
% ------------------------------------------------------------------------------
\begin{table}[htbp]
\centering
\caption{Bảng đặc tả chi tiết Use Case UC07: Xem Báo cáo Thống kê Doanh thu}
\label{tab:uc07_revenue_report}
\vspace{2mm}
\begin{tabularx}{\textwidth}{|l|X|}
\hline
\textbf{Mã Use Case} & \textbf{UC07} \\ \hline
\textbf{Tên Use Case} & Xem Báo cáo Thống kê Doanh thu và Hiệu suất Kinh doanh \\ \hline
\textbf{Tác nhân (Actor)} & Quản trị viên \\ \hline
\textbf{Mô tả tóm tắt} & Quản trị viên truy cập bảng điều khiển Dashboard để xem tổng quan chỉ số kinh doanh, bao gồm doanh thu, số đơn hàng, sản phẩm bán chạy và trạng thái tồn kho. \\ \hline
\textbf{Tiền điều kiện} & Người dùng có quyền \texttt{ROLE\_ADMIN} đã đăng nhập vào hệ thống. \\ \hline
\textbf{Luồng sự kiện chính} & 
\begin{enumerate}[leftmargin=*, nosep]
    \item Quản trị viên truy cập trang Dashboard (\texttt{/dashboard}).
    \item Front-end gửi yêu cầu \texttt{GET /api/v1.0/admin/dashboard} tới Back-end.
    \item Back-end truy vấn cơ sở dữ liệu và tổng hợp các chỉ số kinh doanh (KPI):
    \begin{itemize}
        \item Tổng doanh thu ngày/tuần/tháng hiện tại và tăng trưởng so với kỳ trước.
        \item Tổng số đơn hàng thành công (\texttt{COMPLETED}) trong khoảng thời gian.
        \item Danh sách sản phẩm bán chạy (top items theo số lượng bán).
        \item Số lượng biến thể có tồn kho thấp (cảnh báo sắp hết hàng).
    \end{itemize}
    \item Back-end trả về DTO tổng hợp gồm các trường metric, Front-end hiển thị bảng điều khiển trực quan với các thẻ KPI, biểu đồ và bảng xếp hạng.
\end{enumerate} \\ \hline
\textbf{Luồng ngoại lệ} & 
\begin{itemize}[leftmargin=*, nosep]
    \item \textit{2a. Không có dữ liệu đơn hàng nào trong khoảng thời gian:} Hệ thống hiển thị các thẻ KPI với giá trị bằng $0$ và thông báo chưa có dữ liệu.
\end{itemize} \\ \hline
\textbf{Hậu điều kiện} & Quản trị viên nắm bắt được tình hình kinh doanh tổng quan, doanh thu, sản phẩm bán chạy và trạng thái tồn kho hiện tại. \\ \hline
\end{tabularx}
\end{table}

\subsection{Biểu đồ hoạt động (Activity Diagrams)}

\subsubsection{Biểu đồ hoạt động: Quy trình Bán hàng và Thanh toán tại quầy}
Biểu đồ hoạt động (Hình \ref{fig:activity_checkout}) mô tả tuần tự các bước xử lý nghiệp vụ từ lúc thu ngân chọn mặt hàng, hệ thống đánh giá khuyến mãi tối ưu, đến khi hoàn tất giao dịch thanh toán và trừ kho.

\begin{figure}[H]
\centering
\begin{minipage}{0.9\textwidth}
\centering
\fbox{
    \parbox{0.95\textwidth}{
        \centering \vspace{2mm}
        \textbf{QUY TRÌNH HOẠT ĐỘNG: BÁN HÀNG VÀ THANH TOÁN POS}\\[3mm]
        [Bắt đầu] $\longrightarrow$ [Thu ngân chọn Sản phẩm \& Biến thể SKU]\\[1mm]
        $\downarrow$\\[1mm]
        [Có chọn tùy chọn Modifier?] $\stackrel{Có}{\longrightarrow}$ [Chọn Topping/Đường/Đá] $\longrightarrow$ [Thêm vào Giỏ hàng]\\[1mm]
        $\downarrow$ \textit{Không}\\[1mm]
        [Thêm vào Giỏ hàng] $\longleftarrow$ [Nhận diện Khách hàng thân thiết CRM]\\[1mm]
        $\downarrow$\\[1mm]
        [Hệ thống tự động đánh giá Khuyến mãi tối ưu (No Stacking)]\\[1mm]
        $\downarrow$\\[1mm]
        [Chọn phương thức thanh toán]\\[1mm]
        $\swarrow$ \hspace{4cm} $\searrow$\\[1mm]
        \textbf{[TIỀN MẶT (CASH)]} \hspace{3cm} \textbf{[CHUYỂN KHOẢN VIETQR (PAYOS)]}\\[1mm]
        $\downarrow$ \hspace{5.5cm} $\downarrow$\\[1mm]
        [Nhập tiền khách đưa] \hspace{2.5cm} [Sinh mã VietQR động \& Bật Polling 3s]\\[1mm]
        $\downarrow$ \hspace{5.5cm} $\downarrow$\\[1mm]
        [Tính tiền thừa] \hspace{3cm} [Khách quét mã $\rightarrow$ PayOS gửi Webhook]\\[1mm]
        $\searrow$ \hspace{4cm} $\swarrow$\\[1mm]
        [Xác nhận thanh toán thành công $\rightarrow$ Chuyển trạng thái COMPLETED]\\[1mm]
        $\downarrow$\\[1mm]
        [Ghi giao dịch xuất kho OUT vào sổ cái $\rightarrow$ Cập nhật tồn kho]\\[1mm]
        $\downarrow$\\[1mm]
        [In hóa đơn bán lẻ nhiệt 80mm] $\longrightarrow$ [Kết thúc]\\[2mm]
    }
}
\end{minipage}
\caption{Biểu đồ hoạt động quy trình Bán hàng và Thanh toán tại quầy POS}
\label{fig:activity_checkout}
\end{figure}

\subsubsection{Biểu đồ hoạt động: Quy trình Kiểm kê và Cân bằng tồn kho}
Biểu đồ hoạt động (Hình \ref{fig:activity_stockcheck}) mô tả quy trình kiểm kê định kỳ tại cửa hàng, tính toán sai lệch giữa số lượng thực tế và số liệu trên hệ thống, sau đó thực hiện ghi giao dịch điều chỉnh sổ cái (\texttt{ADJUSTMENT}).

\begin{figure}[H]
\centering
\begin{minipage}{0.9\textwidth}
\centering
\fbox{
    \parbox{0.95\textwidth}{
        \centering \vspace{2mm}
        \textbf{QUY TRÌNH HOẠT ĐỘNG: KIỂM KÊ VÀ CÂN BẰNG TỒN KHO}\\[3mm]
        [Bắt đầu] $\longrightarrow$ [Quản lý mở danh sách Biến thể SKU sản phẩm]\\[1mm]
        $\downarrow$\\[1mm]
        [Hệ thống đọc số dư tồn kho sổ cái hiện tại (\texttt{cachedStockQuantity})]\\[1mm]
        $\downarrow$\\[1mm]
        [Quản lý nhập số lượng đếm thực tế tại quầy]\\[1mm]
        $\downarrow$\\[1mm]
        [Hệ thống tự động tính độ lệch $\Delta = \text{Số thực tế} - \text{Số hệ thống}$]\\[1mm]
        $\downarrow$\\[1mm]
        [Quản lý nhập lý do kiểm kê (hao hụt, nhập sót) \& Nhấn xác nhận]\\[1mm]
        $\downarrow$\\[1mm]
        [Tạo giao dịch ADJUSTMENT với số lượng $\Delta$ vào \texttt{tbl\_inventory\_transactions}]\\[1mm]
        $\downarrow$\\[1mm]
        [Hệ thống tính lại tổng tồn kho sổ cái $\rightarrow$ Cập nhật \texttt{cachedStockQuantity}]\\[1mm]
        $\downarrow$\\[1mm]
        [Ghi vết thao tác vào \texttt{tbl\_activity\_logs}] $\longrightarrow$ [Kết thúc]\\[2mm]
    }
}
\end{minipage}
\caption{Biểu đồ hoạt động quy trình Kiểm kê và Cân bằng tồn kho}
\label{fig:activity_stockcheck}
\end{figure}

\subsubsection{Biểu đồ hoạt động: Quy trình Nhập kho hàng hóa}
Biểu đồ hoạt động (Hình \ref{fig:activity_goods_receipt}) mô tả quy trình nhập kho hàng hóa mới từ nhà cung cấp hoặc bổ sung số lượng cho các biến thể SKU hiện có, ghi nhận giao dịch sổ cái loại \texttt{IN}.

\begin{figure}[H]
\centering
\begin{minipage}{0.9\textwidth}
\centering
\fbox{
    \parbox{0.95\textwidth}{
        \centering \vspace{2mm}
        \textbf{QUY TRÌNH HOẠT ĐỘNG: NHẬP KHO HÀNG HÓA (STOCK IN)}\\[3mm]
        [Bắt đầu] $\longrightarrow$ [Quản trị viên mở trang Quản lý Kho (\texttt{/inventory})]\\[1mm]
        $\downarrow$\\[1mm]
        [Chọn Sản phẩm và Biến thể SKU cần nhập kho]\\[1mm]
        $\downarrow$\\[1mm]
        [Nhập số lượng hàng nhập thêm (quantity $> 0$)]\\[1mm]
        $\downarrow$\\[1mm]
        [Nhập ghi chú nguồn gốc lô hàng (ví dụ: ''Nhập từ NCC ABC'')]\\[1mm]
        $\downarrow$\\[1mm]
        [Nhấn xác nhận nhập kho]\\[1mm]
        $\downarrow$\\[1mm]
        [Back-end tạo giao dịch sổ cái \texttt{IN} trong \texttt{tbl\_inventory\_transactions}]\\[1mm]
        $\downarrow$\\[1mm]
        [Hệ thống cộng thêm số lượng vào tồn kho $\rightarrow$ Cập nhật \texttt{cachedStockQuantity}]\\[1mm]
        $\downarrow$\\[1mm]
        [Ghi vết thao tác vào \texttt{tbl\_activity\_logs}] $\longrightarrow$ [Kết thúc]\\[2mm]
    }
}
\end{minipage}
\caption{Biểu đồ hoạt động quy trình Nhập kho hàng hóa từ nhà cung cấp}
\label{fig:activity_goods_receipt}
\end{figure}

\subsection{Biểu đồ tuần tự (Sequence Diagrams)}

\subsubsection{Biểu đồ tuần tự: Xác thực đăng nhập và Xoay vòng Token (ADR-0007)}
Biểu đồ tuần tự (Hình \ref{fig:seq_auth}) đặc tả luồng tương tác giữa Trình duyệt (Client), Bộ lọc an ninh (JwtRequestFilter), Bộ điều khiển (AuthController), Dịch vụ làm mới mã (RefreshTokenService) và Cơ sở dữ liệu MySQL.

\begin{figure}[H]
\centering
\begin{minipage}{0.95\textwidth}
\centering
\fbox{
    \parbox{0.95\textwidth}{
        \centering \vspace{2mm}
        \textbf{BIỂU ĐỒ TUẦN TỰ: ĐĂNG NHẬP VÀ XOAY VÒNG TOKEN}\\[3mm]
        \begin{tabularx}{\textwidth}{p{2.5cm}cp{3cm}cp{3cm}}
            \textbf{[Client React]} & & \textbf{[Spring Security / Controller]} & & \textbf{[MySQL Database]} \tabularnewline
            $\vert$ & \textit{1. POST /login (email, pass)} $\longrightarrow$ & $\vert$ & & $\vert$ \tabularnewline
            $\vert$ & & $\vert$ $\rightarrow$ Đối chiếu mật khẩu BCrypt & & $\vert$ \tabularnewline
            $\vert$ & & $\vert$ $\rightarrow$ Sinh AccessToken + RefreshToken & & $\vert$ \tabularnewline
            $\vert$ & & $\vert$ \textit{2. Lưu SHA-256 Hash RefreshToken} $\longrightarrow$ & $\vert$ \tabularnewline
            $\vert$ & $\longleftarrow$ \textit{3. Set-Cookie: HttpOnly refreshToken} & $\vert$ & & $\vert$ \tabularnewline
            $\vert$ & $\longleftarrow$ \textit{4. Response JSON: AccessToken} & $\vert$ & & $\vert$ \tabularnewline
            $\vert$ & & $\vert$ & & $\vert$ \tabularnewline
            \multicolumn{5}{c}{\textit{--- Sau 15 phút, Access Token hết hạn $\rightarrow$ Tự động xoay vòng Token ---}} \tabularnewline
            $\vert$ & & $\vert$ & & $\vert$ \tabularnewline
            $\vert$ & \textit{5. POST /auth/refresh (Cookie HttpOnly)} $\longrightarrow$ & $\vert$ & & $\vert$ \tabularnewline
            $\vert$ & & $\vert$ \textit{6. Kiểm tra tokenHash \& familyId} $\longrightarrow$ & $\vert$ \tabularnewline
            $\vert$ & & $\vert$ $\rightarrow$ Thu hồi token cũ (\texttt{revokedAt}) & & $\vert$ \tabularnewline
            $\vert$ & & $\vert$ $\rightarrow$ Sinh cặp Token mới & & $\vert$ \tabularnewline
            $\vert$ & $\longleftarrow$ \textit{7. Set-Cookie mới \& Trả AccessToken mới} & $\vert$ & & $\vert$ \tabularnewline
        \end{tabularx}\\[2mm]
    }
}
\end{minipage}
\caption{Biểu đồ tuần tự quá trình Xác thực đăng nhập và Xoay vòng Token}
\label{fig:seq_auth}
\end{figure}

\subsubsection{Biểu đồ tuần tự: Tạo đơn hàng và Trừ kho tức thời qua Sổ cái}
Biểu đồ tuần tự (Hình \ref{fig:seq_order_create}) thể hiện quá trình phối hợp đồng bộ giữa tầng Front-end, tầng Service nghiệp vụ (\texttt{OrderServiceImpl}, \texttt{PromotionServiceImpl}, \texttt{InventoryServiceImpl}) và cơ sở dữ liệu MySQL khi thực hiện tạo đơn hàng và trừ kho.

\begin{figure}[H]
\centering
\begin{minipage}{0.95\textwidth}
\centering
\fbox{
    \parbox{0.95\textwidth}{
        \centering \vspace{2mm}
        \textbf{BIỂU ĐỒ TUẦN TỰ: TẠO ĐƠN HÀNG VÀ TRỪ KHO SỔ CÁI}\\[3mm]
        \begin{tabularx}{\textwidth}{p{2.2cm}cp{3cm}cp{3.5cm}}
            \textbf{[Client POS]} & & \textbf{[OrderServiceImpl]} & & \textbf{[Promotion / Inventory / DB]} \tabularnewline
            $\vert$ & \textit{1. POST /orders (OrderRequest)} $\longrightarrow$ & $\vert$ & & $\vert$ \tabularnewline
            $\vert$ & & $\vert$ \textit{2. Đánh giá lại Khuyến mãi} $\longrightarrow$ & \textbf{[PromotionService]} \tabularnewline
            $\vert$ & & $\vert$ $\longleftarrow$ \textit{3. Trả về chiết khấu tối ưu} & $\vert$ \tabularnewline
            $\vert$ & & $\vert$ $\rightarrow$ Tính Tổng tiền = Subtotal - Promo + VAT & $\vert$ \tabularnewline
            $\vert$ & & $\vert$ \textit{4. Ghi bản ghi OUT vào sổ cái} $\longrightarrow$ & \textbf{[InventoryTransaction]} \tabularnewline
            $\vert$ & & $\vert$ \textit{5. Cập nhật tồn kho biến thể} $\longrightarrow$ & \textbf{[Variant.cachedStock]} \tabularnewline
            $\vert$ & & $\vert$ \textit{6. Lưu đơn hàng và chi tiết đơn} $\longrightarrow$ & \textbf{[OrderEntity \& OrderItem]} \tabularnewline
            $\vert$ & & $\vert$ \textit{7. Cập nhật tích lũy chi tiêu CRM} $\longrightarrow$ & \textbf{[CustomerEntity]} \tabularnewline
            $\vert$ & $\longleftarrow$ \textit{8. Trả về OrderResponse thành công} & $\vert$ & & $\vert$ \tabularnewline
        \end{tabularx}\\[2mm]
    }
}
\end{minipage}
\caption{Biểu đồ tuần tự quá trình Tạo đơn hàng và Trừ kho sổ cái}
\label{fig:seq_order_create}
\end{figure}

\subsection{Biểu đồ lớp (Class Diagram)}
Cấu trúc quan hệ giữa các lớp đối tượng thực thể (JPA Entities) và các lớp xử lý nghiệp vụ chính được mô tả chi tiết tại Hình \ref{fig:class_diagram}.

\begin{figure}[H]
\centering
\begin{minipage}{0.95\textwidth}
\centering
\fbox{
    \parbox{0.95\textwidth}{
        \centering \vspace{2mm}
        \textbf{SƠ ĐỒ BIỂU ĐỒ LỚP (CLASS DIAGRAM) CÁC ĐỐI TƯỢNG CỐT LÕI}\\[3mm]
        \begin{tabularx}{\textwidth}{|X|c|X|}
            \hline
            \textbf{CategoryEntity} & \multirow{5}{*}{$\stackrel{1}{\longleftrightarrow}\stackrel{N}{}$} & \textbf{ItemEntity} \\
            - id: Long & & - id: Long \\
            - categoryId: String & & - itemId: String \\
            - name: String & & - name: String, price: BigDecimal \\
            - bgColor, imgUrl: String & & - variants: List$<$VariantEntity$>$ \\
            \cline{1-1} \cline{3-3}
            \multicolumn{3}{c}{\vspace{1mm}} \\
            \hline
            \textbf{VariantEntity} & \multirow{5}{*}{$\stackrel{1}{\longleftrightarrow}\stackrel{N}{}$} & \textbf{InventoryTransactionEntity} \\
            - id: Long & & - id: Long \\
            - variantId, sku: String & & - transactionId: String \\
            - basePrice: BigDecimal & & - transactionType: TransactionType \\
            - attributes: Map$<$String, String$>$ & & - quantity: Integer (âm khi xuất) \\
            - cachedStockQuantity: Integer & & - referenceId, note: String \\
            \hline
            \multicolumn{3}{c}{\vspace{1mm}} \\
            \hline
            \textbf{OrderEntity} & \multirow{5}{*}{$\stackrel{1}{\longleftrightarrow}\stackrel{N}{}$} & \textbf{OrderItemEntity} \\
            - id: Long & & - id: Long \\
            - orderId: String & & - itemId, variantId: String \\
            - subtotal, tax, grandTotal: Double & & - name: String, basePrice: Double \\
            - paymentMethod: PaymentMethod & & - quantity: Integer \\
            - paymentDetails: PaymentDetails & & - selectedModifiers: String (JSON) \\
            \hline
        \end{tabularx}\\[2mm]
    }
}
\end{minipage}
\caption{Biểu đồ lớp các thực thể trung tâm của hệ thống}
\label{fig:class_diagram}
\end{figure}

\subsection{Biểu đồ thành phần (Component Diagram)}
Biểu đồ thành phần (Hình \ref{fig:component_diagram}) mô tả các thành phần phần mềm chính của hệ thống POS và mối quan hệ phụ thuộc giữa chúng theo kiến trúc 3 tầng (Presentation - Business Logic - Data Access):

\begin{figure}[H]
\centering
\begin{minipage}{0.95\textwidth}
\centering
\fbox{
    \parbox{0.95\textwidth}{
        \centering \vspace{2mm}
        \textbf{SƠ ĐỒ THÀNH PHẦN HỆ THỐNG POS (COMPONENT DIAGRAM)}\\[3mm]
        \rule{0.9\textwidth}{0.5pt}\\[2mm]
        \fbox{\parbox{0.92\textwidth}{\centering \textbf{TẦNG TRÌNH BÀY (Presentation Layer - React 19 SPA)}\\[2mm]
        \begin{tabularx}{0.9\textwidth}{XXX}
            \fbox{\parbox{2.5cm}{\centering \small \texttt{POS Checkout}\\(Explore, Cart)}} &
            \fbox{\parbox{2.5cm}{\centering \small \texttt{Admin Pages}\\(Dashboard, CRUD)}} &
            \fbox{\parbox{2.5cm}{\centering \small \texttt{Auth Module}\\(Login, Session)}}
        \end{tabularx}\\[1mm]
        \begin{tabularx}{0.9\textwidth}{XX}
            \fbox{\parbox{3.5cm}{\centering \small \texttt{TanStack Query}\\(Server State Cache)}} &
            \fbox{\parbox{3.5cm}{\centering \small \texttt{Axios Interceptor}\\(JWT Bearer, Refresh)}}
        \end{tabularx}
        }}\\[3mm]
        $\Downarrow$ \textit{RESTful JSON (HTTPS)}\\[2mm]
        \fbox{\parbox{0.92\textwidth}{\centering \textbf{TẦNG NGHIỆP VỤ (Business Logic Layer - Spring Boot 4.1.0)}\\[2mm]
        \begin{tabularx}{0.9\textwidth}{XXX}
            \fbox{\parbox{2.5cm}{\centering \small \texttt{AuthController}\\+ JwtFilter}} &
            \fbox{\parbox{2.5cm}{\centering \small \texttt{OrderService}\\+ Compensation}} &
            \fbox{\parbox{2.5cm}{\centering \small \texttt{ItemService}\\+ VariantService}}}
        \end{tabularx}\\[1mm]
        \begin{tabularx}{0.9\textwidth}{XXX}
            \fbox{\parbox{2.5cm}{\centering \small \texttt{InventoryService}\\(Ledger IN/OUT)}} &
            \fbox{\parbox{2.5cm}{\centering \small \texttt{PromotionService}\\(No-Stacking)}} &
            \fbox{\parbox{2.5cm}{\centering \small \texttt{ActivityLog}\\(Audit Trail)}}
        \end{tabularx}
        }}\\[3mm]
        $\Downarrow$ \textit{JPA / Flyway}\\[2mm]
        \fbox{\parbox{0.92\textwidth}{\centering \textbf{TẦNG TRUY CẬP DỮ LIỆU (Data Access Layer)}\\[2mm]
        \begin{tabularx}{0.9\textwidth}{XXX}
            \fbox{\parbox{2.5cm}{\centering \small \texttt{MySQL 8.0}\\(14 Tables InnoDB)}} &
            \fbox{\parbox{2.5cm}{\centering \small \texttt{AWS S3 SDK}\\(Image Storage)}} &
            \fbox{\parbox{2.5cm}{\centering \small \texttt{PayOS SDK}\\(VietQR Webhook)}}
        \end{tabularx}
        }}\\[2mm]
    }
}
\end{minipage}
\caption{Biểu đồ thành phần kiến trúc 3 tầng của hệ thống POS}
\label{fig:component_diagram}
\end{figure}

Các thành phần chính được phân tách theo nguyên tắc trách nhiệm đơn lẻ (Single Responsibility):
\begin{itemize}[leftmargin=1.5cm]
    \item \textbf{Tầng trình bày (Presentation):} Giao diện React 19 SPA gồm module POS bán hàng, trang quản trị CRUD và module xác thực. Quản lý trạng thái server qua TanStack Query v5, gắn JWT vào mọi request qua Axios Interceptor.
    \item \textbf{Tầng nghiệp vụ (Business Logic):} Các Service xử lý toàn bộ quy tắc kinh doanh --- tính toán khuyến mãi No-Stacking, giao dịch bù trừ kho khi hủy đơn (Compensating Transaction), kiểm soát truy cập RBAC qua Spring Security và ghi nhật ký kiểm toán thống nhất.
    \item \textbf{Tầng truy cập dữ liệu (Data Access):} MySQL 8.0 InnoDB (14 bảng) quản lý qua Flyway migration, tích hợp AWS S3 cho lưu trữ hình ảnh và PayOS SDK cho cổng thanh toán VietQR.
\end{itemize}

\subsection{Biểu đồ triển khai (Deployment Diagram)}
Kiến trúc phần cứng và các thành phần phần mềm triển khai trên môi trường thực tế được thiết kế theo mô hình phân tán nhiều tầng (Multi-tier Distributed Deployment Architecture), như trình bày tại Hình \ref{fig:deployment_diagram}:

\begin{figure}[H]
\centering
\begin{minipage}{0.95\textwidth}
\centering
\fbox{
    \parbox{0.95\textwidth}{
        \centering \vspace{2mm}
        \textbf{SƠ ĐỒ TRIỂN KHAI HỆ THỐNG THỰC TẾ (DEPLOYMENT DIAGRAM)}\\[3mm]
        \begin{tabularx}{\textwidth}{ccc}
            \fbox{\parbox{3.5cm}{\centering \textbf{THIẾT BỊ ĐẦU CUỐI}\\
            \textit{(Client POS Node)}\\
            $\bullet$ Trình duyệt Web\\
            $\bullet$ React 19 SPA\\
            $\bullet$ Máy in nhiệt 80mm}} 
            & $\stackrel{\text{HTTPS / WSS}}{\longleftrightarrow}$ & 
            \fbox{\parbox{4.5cm}{\centering \textbf{MÁY CHỦ ỨNG DỤNG}\\
            \textit{(Docker Host Linux)}\\
            $\bullet$ \textbf{Nginx Container (Port 80/443)}\\
            $\bullet$ \textbf{Spring Boot Container (Port 8080)}\\
            \hspace{2mm}- Java 25 Runtime Engine\\
            \hspace{2mm}- 13 REST Controllers / 11 Services}} \\[5mm]
            & & $\updownarrow$ \textit{Mạng nội bộ Bridge (Port 3306)}\\[2mm]
            \fbox{\parbox{3.5cm}{\centering \textbf{CỔNG THANH TOÁN}\\
            \textit{(PayOS Cloud API)}\\
            $\bullet$ Sinh mã VietQR động\\
            $\bullet$ Gửi Webhook HMAC}}
            & $\stackrel{\text{HTTPS REST}}{\longleftrightarrow}$ &
            \fbox{\parbox{4.5cm}{\centering \textbf{CƠ SỞ DỮ LIỆU MYSQL}\\
            \textit{(MySQL 8.0 InnoDB Container)}\\
            $\bullet$ Bảng dữ liệu quan hệ (14 tables)\\
            $\bullet$ Sổ cái biến động kho\\
            $\bullet$ Flyway Schema History}} \\[5mm]
            & & $\updownarrow$ \textit{AWS SDK HTTPS}\\[2mm]
            & & \fbox{\parbox{4.5cm}{\centering \textbf{KHO LƯU TRỮ ĐÁM MÂY}\\
            \textit{(Amazon S3 Bucket)}\\
            $\bullet$ Tệp ảnh sản phẩm \& danh mục}}
        \end{tabularx}\\[2mm]
    }
}
\end{minipage}
\caption{Biểu đồ triển khai kiến trúc hệ thống POS trên hạ tầng đám mây}
\label{fig:deployment_diagram}
\end{figure}

Các nút thành phần trong sơ đồ triển khai:
\begin{itemize}[leftmargin=1.5cm]
    \item \textbf{Thiết bị đầu cuối thu ngân (Client Node):} Máy tính hoặc máy tính bảng kết nối máy in nhiệt 80mm qua cổng USB/LAN, chạy ứng dụng ReactJS SPA trên nền trình duyệt hiện đại.
    \item \textbf{Máy chủ Web và Reverse Proxy (Nginx Container):} Đóng vai trò là cổng tiếp nhận toàn bộ các kết nối từ bên ngoài, phục vụ các tệp tĩnh nén Gzip của Frontend và chuyển tiếp (Reverse Proxy) các request có tiền tố \texttt{/api/v1.0} vào container Spring Boot.
    \item \textbf{Máy chủ ứng dụng xử lý nghiệp vụ (Spring Boot Container):} Vận hành trên môi trường Java 25, đảm nhận toàn bộ tính toán nghiệp vụ, kiểm tra an ninh JWT, kết nối cơ sở dữ liệu và liên lạc với các dịch vụ bên ngoài.
    \item \textbf{Máy chủ Cơ sở dữ liệu (MySQL Database Container):} Lưu trữ toàn bộ dữ liệu quan hệ, sổ cái tồn kho và lịch sử kiểm toán với ổ lưu trữ gắn ngoài (Docker Volume Persistent Storage) đảm bảo an toàn dữ liệu.
    \item \textbf{Dịch vụ đám mây tích hợp (PayOS \& AWS S3):} Cung cấp hạ tầng thanh toán ngân hàng tự động và lưu trữ tệp hình ảnh không giới hạn băng thông.
\end{itemize}

% ==============================================================================
% MỤC 2.3: THIẾT KẾ CƠ SỞ DỮ LIỆU
% ==============================================================================
% ==============================================================================
% 02_CHAPTER2_PART2_DATABASE.TEX - THIẾT KẾ CƠ SỞ DỮ LIỆU & TỪ ĐIỂN DỮ LIỆU
% ==============================================================================

\section{Phân tích và thiết kế cơ sở dữ liệu}
\label{sec:thiet_ke_csdl}

\subsection{Mô hình thực thể liên kết (ERD)}

\subsubsection{Mô hình ERD mức quan niệm (Conceptual Data Model)}
Mô hình ERD mức quan niệm xác định các thực thể thế giới thực trong bài toán kinh doanh POS và các mối liên kết bản số giữa chúng:
\begin{itemize}[leftmargin=1.5cm]
    \item \textbf{Danh mục (Category) và Sản phẩm (Item):} Quan hệ $1-N$. Một danh mục (ví dụ: ''Đồ uống'', ''Thức ăn nhanh'') có thể chứa nhiều sản phẩm; mỗi sản phẩm bắt buộc phải thuộc về đúng một danh mục xác định.
    \item \textbf{Sản phẩm (Item) và Biến thể SKU (Variant):} Quan hệ $1-N$. Một sản phẩm cơ sở (như ''Áo thun'') có thể có nhiều biến thể vật lý (SKU) khác nhau về kích cỡ (S/M/L) và màu sắc (Đỏ/Xanh). Biến thể là đơn vị duy nhất quản lý tồn kho và có đơn giá bán riêng biệt.
    \item \textbf{Sản phẩm (Item) và Nhóm tùy chọn (ModifierGroup):} Quan hệ $N-N$. Một sản phẩm có thể liên kết với nhiều nhóm tùy chọn (như ''Mức đường'', ''Topping thêm''); và một nhóm tùy chọn có thể được tái sử dụng cho nhiều sản phẩm khác nhau.
    \item \textbf{Nhóm tùy chọn (ModifierGroup) và Tùy chọn con (Modifier):} Quan hệ $1-N$. Mỗi nhóm tùy chọn chứa danh sách các tùy chọn cụ thể (như Trân châu trắng, Thạch nha đam) kèm mức điều chỉnh giá tương ứng (\texttt{priceAdjustment}).
    \item \textbf{Biến thể (Variant) và Giao dịch sổ cái kho (InventoryTransaction):} Quan hệ $1-N$. Mọi biến động tăng hoặc giảm tồn kho của một biến thể đều được ghi nhận bằng một dòng giao dịch sổ cái độc lập (IN, OUT, ADJUSTMENT).
    \item \textbf{Khách hàng (Customer) và Đơn hàng (Order):} Quan hệ $1-N$. Một khách hàng thân thiết có thể thực hiện nhiều đơn mua hàng; mỗi đơn hàng có thể gắn với một khách hàng hoặc để trống (khách vãng lai).
    \item \textbf{Đơn hàng (Order) và Chi tiết mặt hàng (OrderItem):} Quan hệ $1-N$. Một đơn hàng có thể bao gồm nhiều dòng sản phẩm; mỗi dòng lưu vết chính xác số lượng, đơn giá và cấu hình tùy chọn tại thời điểm giao dịch.
    \item \textbf{Khuyến mãi (Promotion) và Đơn hàng (Order):} Quan hệ $1-N$. Mỗi đơn hàng chỉ được gắn tối đa với một chương trình khuyến mãi duy nhất (No Stacking).
    \item \textbf{Người dùng (User) và Nhật ký kiểm toán (ActivityLog):} Quan hệ $1-N$. Một tài khoản nhân viên hoặc quản trị viên có thể thực hiện nhiều thao tác được ghi log trên hệ thống.
\end{itemize}

\subsubsection{Mô hình ERD mức logic và Sơ đồ quan hệ thực thể}
Trong mô hình mức logic, quan hệ nhiều - nhiều ($N-N$) giữa bảng \texttt{tbl\_items} và \texttt{tbl\_modifier\_groups} được giải quyết triệt để thông qua bảng trung gian \texttt{tbl\_item\_modifier\_groups} chứa cặp khóa ngoại (\texttt{item\_id}, \texttt{modifier\_group\_id}). Đồng thời, các trường quan hệ đều được thiết lập ràng buộc toàn vẹn khóa ngoại (Foreign Key Constraints) với quy tắc xử lý xóa thích hợp (\texttt{RESTRICT} đối với danh mục có sản phẩm, \texttt{CASCADE} đối với sản phẩm cha bị xóa).

Sơ đồ mô hình quan hệ thực thể (ERD) toàn diện của hệ thống được minh họa trực quan tại Hình \ref{fig:erd_diagram}.

\begin{figure}[htbp]
\centering
\begin{minipage}{0.95\textwidth}
\centering
\fbox{
    \parbox{0.95\textwidth}{
        \centering \vspace{2mm}
        {\fontsize{12pt}{14pt}\selectfont \textbf{SƠ ĐỒ QUAN HỆ THỰC THỂ CƠ SỞ DỮ LIỆU (ERD - 13 BẢNG)}}\\[3mm]
        \begin{tabularx}{\textwidth}{ccc}
            \fbox{\parbox{3.5cm}{\centering \textbf{tbl\_category}\\
            \underline{id} (PK)\\
            categoryId (UK)\\
            name, bgColor, imgUrl}} 
            & $\stackrel{1}{\longleftrightarrow}\stackrel{N}{\longrightarrow}$ & 
            \fbox{\parbox{4.5cm}{\centering \textbf{tbl\_items}\\
            \underline{id} (PK), itemId (UK)\\
            name, price, imgUrl\\
            \textbf{category\_id} (FK)}} \\[8mm]
            & & $\updownarrow$ $\stackrel{1}{\longleftrightarrow}\stackrel{N}{}$ \\[2mm]
            \fbox{\parbox{3.5cm}{\centering \textbf{tbl\_item\_modifier\_groups}\\
            \underline{item\_id} (FK, PK)\\
            \underline{modifier\_group\_id} (FK, PK)}}
            & $\stackrel{N}{\longleftarrow}\stackrel{1}{\longleftrightarrow}$ &
            \fbox{\parbox{4.5cm}{\centering \textbf{tbl\_variants}\\
            \underline{id} (PK), variantId (UK)\\
            sku (UK), basePrice\\
            attributes (JSON), cachedStock\\
            \textbf{item\_id} (FK)}} \\[8mm]
            $\updownarrow$ $\stackrel{N}{\longleftarrow}\stackrel{1}{}$ & & $\updownarrow$ $\stackrel{1}{\longleftrightarrow}\stackrel{N}{}$ \\[2mm]
            \fbox{\parbox{3.5cm}{\centering \textbf{tbl\_modifier\_groups}\\
            \underline{id} (PK), groupId (UK)\\
            name, min/maxSelections}}
            & $\stackrel{1}{\longleftrightarrow}\stackrel{N}{\longrightarrow}$ &
            \fbox{\parbox{4.5cm}{\centering \textbf{tbl\_inventory\_transactions}\\
            \underline{id} (PK), transactionId (UK)\\
            transactionType (IN/OUT/ADJ)\\
            quantity, referenceId, note\\
            \textbf{variant\_id} (FK)}} \\[8mm]
            $\updownarrow$ $\stackrel{1}{\longleftrightarrow}\stackrel{N}{}$ & & \\[2mm]
            \fbox{\parbox{3.5cm}{\centering \textbf{tbl\_modifiers}\\
            \underline{id} (PK), modifierId (UK)\\
            name, priceAdjustment\\
            \textbf{group\_id} (FK)}}
            & &
            \fbox{\parbox{4.5cm}{\centering \textbf{tbl\_orders}\\
            \underline{id} (PK - PayOS orderCode)\\
            order\_code (UK), customerId\\
            subtotal, discountAmount, tax\\
            grandTotal, paymentMethod\\
            paymentDetails (Embedded)}} \\[8mm]
            & & $\updownarrow$ $\stackrel{1}{\longleftrightarrow}\stackrel{N}{}$ \\[2mm]
            \fbox{\parbox{3.5cm}{\centering \textbf{tbl\_promotions}\\
            \underline{id} (PK), promotionId (UK)\\
            name, type, code (UK)\\
            discountType, discountValue\\
            startDate, endDate, isActive}}
            & &
            \fbox{\parbox{4.5cm}{\centering \textbf{tbl\_order\_items}\\
            \underline{id} (PK)\\
            itemId, variantId, name\\
            basePrice, price, quantity\\
            selectedModifiers (JSON)\\
            \textbf{order\_id} (FK)}} \\[8mm]
            \fbox{\parbox{3.5cm}{\centering \textbf{tbl\_customers}\\
            \underline{id} (PK), customerId (UK)\\
            phoneNumber (UK), name\\
            totalSpent, orderCount}}
            & &
            \fbox{\parbox{4.5cm}{\centering \textbf{tbl\_users} / \textbf{refresh\_tokens}\\
            \underline{id} (PK), userId / tokenHash (UK)\\
            email, password, role\\
            familyId, issuedAt, revokedAt}} \\[2mm]
        \end{tabularx}\\[2mm]
    }
}
\end{minipage}
\caption{Sơ đồ quan hệ thực thể (ERD) mức logic của hệ thống POS}
\label{fig:erd_diagram}
\end{figure}

\subsection{Từ điển dữ liệu chi tiết (Data Dictionary - 13 bảng)}
Dưới đây là đặc tả chi tiết toàn bộ các trường dữ liệu, kiểu dữ liệu MySQL, ràng buộc toàn vẹn và ý nghĩa nghiệp vụ của 13 bảng thực tế trong cơ sở dữ liệu:

% ------------------------------------------------------------------------------
% 1. BẢNG TBL_USERS
% ------------------------------------------------------------------------------
\subsubsection{Bảng \texttt{tbl\_users} (Tài khoản người dùng và Phân quyền)}
Lưu trữ thông tin tài khoản nhân viên thu ngân và quản trị viên hệ thống (Bảng \ref{tab:tbl_users}).

\begin{table}[htbp]
\centering
\caption{Từ điển dữ liệu bảng \texttt{tbl\_users}}
\label{tab:tbl_users}
\vspace{2mm}
\begin{tabularx}{\textwidth}{|p{2.5cm}|p{2.8cm}|c|c|p{1.8cm}|X|}
\hline
\textbf{Tên trường} & \textbf{Kiểu dữ liệu} & \textbf{Khóa} & \textbf{Not Null} & \textbf{Mặc định} & \textbf{Diễn giải nghiệp vụ} \\ \hline
\texttt{id} & \texttt{BIGINT} & PK & Có & Auto Inc & Khóa chính định danh tự tăng. \\ \hline
\texttt{userId} & \texttt{VARCHAR(255)} & UK & Có & Không & Mã định danh công khai (UUID). \\ \hline
\texttt{email} & \texttt{VARCHAR(255)} & UK & Có & Không & Email đăng nhập duy nhất. \\ \hline
\texttt{password} & \texttt{VARCHAR(255)} & - & Có & Không & Mật khẩu đã băm bằng BCrypt. \\ \hline
\texttt{role} & \texttt{VARCHAR(50)} & - & Có & Không & Vai trò: \texttt{ROLE\_ADMIN} hoặc \texttt{ROLE\_USER}. \\ \hline
\texttt{name} & \texttt{VARCHAR(255)} & - & Có & Không & Họ và tên đầy đủ của nhân viên. \\ \hline
\texttt{createdAt} & \texttt{DATETIME(6)} & - & Có & CURRENT & Thời điểm khởi tạo tài khoản. \\ \hline
\texttt{updatedAt} & \texttt{DATETIME(6)} & - & Không & NULL & Thời điểm cập nhật gần nhất. \\ \hline
\end{tabularx}
\end{table}

% ------------------------------------------------------------------------------
% 2. BẢNG REFRESH_TOKENS
% ------------------------------------------------------------------------------
\subsubsection{Bảng \texttt{refresh\_tokens} (Quản lý phiên làm việc bảo mật)}
Lưu vết các khóa làm mới dài hạn phục vụ cơ chế Token Rotation chống Replay Attack (Bảng \ref{tab:refresh_tokens}).

\begin{table}[htbp]
\centering
\caption{Từ điển dữ liệu bảng \texttt{refresh\_tokens}}
\label{tab:refresh_tokens}
\vspace{2mm}
\begin{tabularx}{\textwidth}{|p{2.5cm}|p{2.8cm}|c|c|p{1.8cm}|X|}
\hline
\textbf{Tên trường} & \textbf{Kiểu dữ liệu} & \textbf{Khóa} & \textbf{Not Null} & \textbf{Mặc định} & \textbf{Diễn giải nghiệp vụ} \\ \hline
\texttt{id} & \texttt{BIGINT} & PK & Có & Auto Inc & Khóa chính định danh tự tăng. \\ \hline
\texttt{tokenHash} & \texttt{VARCHAR(128)} & UK & Có & Không & Mã băm SHA-256 của Refresh Token. \\ \hline
\texttt{userEmail} & \texttt{VARCHAR(255)} & Index & Có & Không & Email người dùng sở hữu phiên. \\ \hline
\texttt{familyId} & \texttt{VARCHAR(36)} & Index & Có & Không & Mã định danh nhóm phiên xoay vòng. \\ \hline
\texttt{issuedAt} & \texttt{DATETIME(6)} & - & Có & Không & Thời điểm cấp phát token. \\ \hline
\texttt{expiresAt} & \texttt{DATETIME(6)} & - & Có & Không & Thời điểm hết hạn (7 ngày). \\ \hline
\texttt{revokedAt} & \texttt{DATETIME(6)} & - & Không & NULL & Thời điểm thu hồi phiên (nếu bị hủy). \\ \hline
\end{tabularx}
\end{table}

% ------------------------------------------------------------------------------
% 3. BẢNG TBL_CATEGORY
% ------------------------------------------------------------------------------
\subsubsection{Bảng \texttt{tbl\_category} (Danh mục sản phẩm)}
Lưu trữ các nhóm danh mục phân loại mặt hàng (Bảng \ref{tab:tbl_category}).

\begin{table}[htbp]
\centering
\caption{Từ điển dữ liệu bảng \texttt{tbl\_category}}
\label{tab:tbl_category}
\vspace{2mm}
\begin{tabularx}{\textwidth}{|p{2.5cm}|p{2.8cm}|c|c|p{1.8cm}|X|}
\hline
\textbf{Tên trường} & \textbf{Kiểu dữ liệu} & \textbf{Khóa} & \textbf{Not Null} & \textbf{Mặc định} & \textbf{Diễn giải nghiệp vụ} \\ \hline
\texttt{id} & \texttt{BIGINT} & PK & Có & Auto Inc & Khóa chính định danh tự tăng. \\ \hline
\texttt{categoryId} & \texttt{VARCHAR(255)} & UK & Có & Không & Mã định danh danh mục (UUID). \\ \hline
\texttt{name} & \texttt{VARCHAR(255)} & UK & Có & Không & Tên danh mục (ví dụ: Nước uống). \\ \hline
\texttt{description} & \texttt{TEXT} & - & Không & NULL & Mô tả chi tiết danh mục. \\ \hline
\texttt{bgColor} & \texttt{VARCHAR(50)} & - & Không & NULL & Mã màu nền hiển thị trên giao diện POS. \\ \hline
\texttt{imgUrl} & \texttt{VARCHAR(1000)} & - & Không & NULL & Đường dẫn ảnh đại diện trên AWS S3. \\ \hline
\texttt{createdAt} & \texttt{DATETIME(6)} & - & Có & CURRENT & Thời điểm tạo danh mục. \\ \hline
\texttt{updatedAt} & \texttt{DATETIME(6)} & - & Không & NULL & Thời điểm cập nhật gần nhất. \\ \hline
\end{tabularx}
\end{table}

% ------------------------------------------------------------------------------
% 4. BẢNG TBL_ITEMS
% ------------------------------------------------------------------------------
\subsubsection{Bảng \texttt{tbl\_items} (Mặt hàng / Sản phẩm cơ sở)}
Lưu trữ thông tin sản phẩm cơ sở trước khi phân tách thành các biến thể (Bảng \ref{tab:tbl_items}).

\begin{table}[htbp]
\centering
\caption{Từ điển dữ liệu bảng \texttt{tbl\_items}}
\label{tab:tbl_items}
\vspace{2mm}
\begin{tabularx}{\textwidth}{|p{2.5cm}|p{2.8cm}|c|c|p{1.8cm}|X|}
\hline
\textbf{Tên trường} & \textbf{Kiểu dữ liệu} & \textbf{Khóa} & \textbf{Not Null} & \textbf{Mặc định} & \textbf{Diễn giải nghiệp vụ} \\ \hline
\texttt{id} & \texttt{BIGINT} & PK & Có & Auto Inc & Khóa chính định danh tự tăng. \\ \hline
\texttt{itemId} & \texttt{VARCHAR(255)} & UK & Có & Không & Mã sản phẩm công khai (UUID). \\ \hline
\texttt{name} & \texttt{VARCHAR(255)} & - & Có & Không & Tên mặt hàng cơ sở (ví dụ: Trà sữa). \\ \hline
\texttt{description} & \texttt{TEXT} & - & Không & NULL & Mô tả thông tin sản phẩm. \\ \hline
\texttt{price} & \texttt{DECIMAL(19,2)} & - & Có & Không & Đơn giá bán cơ sở mặc định. \\ \hline
\texttt{imgUrl} & \texttt{VARCHAR(1000)} & - & Không & NULL & Đường dẫn ảnh sản phẩm trên AWS S3. \\ \hline
\texttt{category\_id} & \texttt{BIGINT} & FK & Có & Không & Liên kết tới \texttt{tbl\_category.id} (RESTRICT). \\ \hline
\texttt{createdAt} & \texttt{DATETIME(6)} & - & Có & CURRENT & Thời điểm tạo sản phẩm. \\ \hline
\texttt{updatedAt} & \texttt{DATETIME(6)} & - & Không & NULL & Thời điểm cập nhật sản phẩm. \\ \hline
\end{tabularx}
\end{table}

% ------------------------------------------------------------------------------
% 5. BẢNG TBL_VARIANTS
% ------------------------------------------------------------------------------
\subsubsection{Bảng \texttt{tbl\_variants} (Biến thể sản phẩm SKU \& Tồn kho)}
Quản lý các biến thể vật lý SKU, lưu trữ thuộc tính động và số dư tồn kho đệm (Bảng \ref{tab:tbl_variants}).

\begin{table}[htbp]
\centering
\caption{Từ điển dữ liệu bảng \texttt{tbl\_variants}}
\label{tab:tbl_variants}
\vspace{2mm}
\begin{tabularx}{\textwidth}{|p{2.5cm}|p{2.8cm}|c|c|p{1.8cm}|X|}
\hline
\textbf{Tên trường} & \textbf{Kiểu dữ liệu} & \textbf{Khóa} & \textbf{Not Null} & \textbf{Mặc định} & \textbf{Diễn giải nghiệp vụ} \\ \hline
\texttt{id} & \texttt{BIGINT} & PK & Có & Auto Inc & Khóa chính định danh tự tăng. \\ \hline
\texttt{variantId} & \texttt{VARCHAR(255)} & UK & Có & Không & Mã định danh biến thể (UUID). \\ \hline
\texttt{sku} & \texttt{VARCHAR(255)} & UK & Có & Không & Mã vạch quản lý tồn kho duy nhất. \\ \hline
\texttt{basePrice} & \texttt{DECIMAL(19,2)} & - & Có & Không & Giá bán áp dụng cho biến thể này. \\ \hline
\texttt{attributes} & \texttt{TEXT} & - & Không & NULL & Chuỗi JSON thuộc tính (Size, Color). \\ \hline
\texttt{cachedStock} & \texttt{INT} & - & Có & 0 & Số lượng tồn kho đệm đồng bộ từ sổ cái. \\ \hline
\texttt{item\_id} & \texttt{BIGINT} & FK & Có & Không & Khóa ngoại tới \texttt{tbl\_items.id} (CASCADE). \\ \hline
\texttt{createdAt} & \texttt{DATETIME(6)} & - & Có & CURRENT & Thời điểm khởi tạo biến thể. \\ \hline
\texttt{updatedAt} & \texttt{DATETIME(6)} & - & Không & NULL & Thời điểm cập nhật biến thể. \\ \hline
\end{tabularx}
\end{table}

% ------------------------------------------------------------------------------
% 6. BẢNG TBL_MODIFIER_GROUPS
% ------------------------------------------------------------------------------
\subsubsection{Bảng \texttt{tbl\_modifier\_groups} (Nhóm tùy chọn bổ sung)}
Quản lý các nhóm tùy chọn chế biến hoặc gia giảm (Bảng \ref{tab:tbl_modifier_groups}).

\begin{table}[htbp]
\centering
\caption{Từ điển dữ liệu bảng \texttt{tbl\_modifier\_groups}}
\label{tab:tbl_modifier_groups}
\vspace{2mm}
\begin{tabularx}{\textwidth}{|p{2.5cm}|p{2.8cm}|c|c|p{1.8cm}|X|}
\hline
\textbf{Tên trường} & \textbf{Kiểu dữ liệu} & \textbf{Khóa} & \textbf{Not Null} & \textbf{Mặc định} & \textbf{Diễn giải nghiệp vụ} \\ \hline
\texttt{id} & \texttt{BIGINT} & PK & Có & Auto Inc & Khóa chính định danh tự tăng. \\ \hline
\texttt{groupId} & \texttt{VARCHAR(255)} & UK & Có & Không & Mã định danh nhóm tùy chọn (UUID). \\ \hline
\texttt{name} & \texttt{VARCHAR(255)} & - & Có & Không & Tên nhóm (ví dụ: Mức đường, Topping). \\ \hline
\texttt{description} & \texttt{TEXT} & - & Không & NULL & Mô tả thông tin nhóm tùy chọn. \\ \hline
\texttt{minSelections} & \texttt{INT} & - & Có & 0 & Số lượng chọn tối thiểu bắt buộc. \\ \hline
\texttt{maxSelections} & \texttt{INT} & - & Có & 1 & Số lượng chọn tối đa cho phép. \\ \hline
\texttt{createdAt} & \texttt{DATETIME(6)} & - & Có & CURRENT & Thời điểm khởi tạo nhóm. \\ \hline
\texttt{updatedAt} & \texttt{DATETIME(6)} & - & Không & NULL & Thời điểm cập nhật nhóm. \\ \hline
\end{tabularx}
\end{table}

% ------------------------------------------------------------------------------
% 7. BẢNG TBL_MODIFIERS
% ------------------------------------------------------------------------------
\subsubsection{Bảng \texttt{tbl\_modifiers} (Tùy chọn bổ sung cụ thể)}
Lưu trữ các lựa chọn chi tiết và mức phụ thu tương ứng (Bảng \ref{tab:tbl_modifiers}).

\begin{table}[htbp]
\centering
\caption{Từ điển dữ liệu bảng \texttt{tbl\_modifiers}}
\label{tab:tbl_modifiers}
\vspace{2mm}
\begin{tabularx}{\textwidth}{|p{2.5cm}|p{2.8cm}|c|c|p{1.8cm}|X|}
\hline
\textbf{Tên trường} & \textbf{Kiểu dữ liệu} & \textbf{Khóa} & \textbf{Not Null} & \textbf{Mặc định} & \textbf{Diễn giải nghiệp vụ} \\ \hline
\texttt{id} & \texttt{BIGINT} & PK & Có & Auto Inc & Khóa chính định danh tự tăng. \\ \hline
\texttt{modifierId} & \texttt{VARCHAR(255)} & UK & Có & Không & Mã định danh tùy chọn (UUID). \\ \hline
\texttt{name} & \texttt{VARCHAR(255)} & - & Có & Không & Tên tùy chọn (ví dụ: Trân châu trắng). \\ \hline
\texttt{priceAdj} & \texttt{DECIMAL(19,2)} & - & Có & 0.00 & Mức phụ thu cộng thêm vào giá bán. \\ \hline
\texttt{group\_id} & \texttt{BIGINT} & FK & Có & Không & Khóa ngoại tới \texttt{tbl\_modifier\_groups.id}. \\ \hline
\texttt{createdAt} & \texttt{DATETIME(6)} & - & Có & CURRENT & Thời điểm tạo tùy chọn. \\ \hline
\texttt{updatedAt} & \texttt{DATETIME(6)} & - & Không & NULL & Thời điểm cập nhật tùy chọn. \\ \hline
\end{tabularx}
\end{table}

% ------------------------------------------------------------------------------
% 8. BẢNG TBL_ITEM_MODIFIER_GROUPS
% ------------------------------------------------------------------------------
\subsubsection{Bảng \texttt{tbl\_item\_modifier\_groups} (Bảng trung gian Sản phẩm - Tùy chọn)}
Phân giải mối quan hệ nhiều - nhiều ($N-N$) giữa sản phẩm và nhóm tùy chọn (Bảng \ref{tab:tbl_item_mod}).

\begin{table}[htbp]
\centering
\caption{Từ điển dữ liệu bảng \texttt{tbl\_item\_modifier\_groups}}
\label{tab:tbl_item_mod}
\vspace{2mm}
\begin{tabularx}{\textwidth}{|p{3.5cm}|p{2.8cm}|c|c|X|}
\hline
\textbf{Tên trường} & \textbf{Kiểu dữ liệu} & \textbf{Khóa} & \textbf{Not Null} & \textbf{Diễn giải nghiệp vụ} \\ \hline
\texttt{item\_id} & \texttt{BIGINT} & PK, FK & Có & Khóa ngoại liên kết tới \texttt{tbl\_items.id}. \\ \hline
\texttt{modifier\_group\_id} & \texttt{BIGINT} & PK, FK & Có & Khóa ngoại liên kết tới \texttt{tbl\_modifier\_groups.id}. \\ \hline
\end{tabularx}
\end{table}

% ------------------------------------------------------------------------------
% 9. BẢNG TBL_INVENTORY_TRANSACTIONS
% ------------------------------------------------------------------------------
\subsubsection{Bảng \texttt{tbl\_inventory\_transactions} (Sổ cái giao dịch kho)}
Lưu vết toàn bộ biến động kho vật lý theo nguyên tắc kế toán kép (Bảng \ref{tab:tbl_inventory}).

\begin{table}[htbp]
\centering
\caption{Từ điển dữ liệu bảng \texttt{tbl\_inventory\_transactions}}
\label{tab:tbl_inventory}
\vspace{2mm}
\begin{tabularx}{\textwidth}{|p{2.5cm}|p{2.8cm}|c|c|p{1.8cm}|X|}
\hline
\textbf{Tên trường} & \textbf{Kiểu dữ liệu} & \textbf{Khóa} & \textbf{Not Null} & \textbf{Mặc định} & \textbf{Diễn giải nghiệp vụ} \\ \hline
\texttt{id} & \texttt{BIGINT} & PK & Có & Auto Inc & Khóa chính định danh tự tăng. \\ \hline
\texttt{transactionId} & \texttt{VARCHAR(255)} & UK & Có & Không & Mã giao dịch sổ cái (UUID). \\ \hline
\texttt{variant\_id} & \texttt{BIGINT} & FK & Có & Không & Khóa ngoại liên kết \texttt{tbl\_variants.id}. \\ \hline
\texttt{transType} & \texttt{VARCHAR(20)} & - & Có & Không & Loại: \texttt{IN} (Nhập), \texttt{OUT} (Xuất), \texttt{ADJUSTMENT} (Cân bằng). \\ \hline
\texttt{quantity} & \texttt{INT} & - & Có & Không & Số lượng biến động (âm khi xuất). \\ \hline
\texttt{referenceId} & \texttt{VARCHAR(255)} & - & Không & NULL & Mã đơn hàng hoặc mã phiếu kiểm kho. \\ \hline
\texttt{note} & \texttt{TEXT} & - & Không & NULL & Ghi chú nguyên nhân biến động kho. \\ \hline
\texttt{createdAt} & \texttt{DATETIME(6)} & - & Có & CURRENT & Thời điểm phát sinh giao dịch. \\ \hline
\end{tabularx}
\end{table}

% ------------------------------------------------------------------------------
% 10. BẢNG TBL_CUSTOMERS
% ------------------------------------------------------------------------------
\subsubsection{Bảng \texttt{tbl\_customers} (Khách hàng thân thiết CRM)}
Quản lý thông tin và tích lũy doanh số của khách hàng thân thiết (Bảng \ref{tab:tbl_customers}).

\begin{table}[htbp]
\centering
\caption{Từ điển dữ liệu bảng \texttt{tbl\_customers}}
\label{tab:tbl_customers}
\vspace{2mm}
\begin{tabularx}{\textwidth}{|p{2.5cm}|p{2.8cm}|c|c|p{1.8cm}|X|}
\hline
\textbf{Tên trường} & \textbf{Kiểu dữ liệu} & \textbf{Khóa} & \textbf{Not Null} & \textbf{Mặc định} & \textbf{Diễn giải nghiệp vụ} \\ \hline
\texttt{id} & \texttt{BIGINT} & PK & Có & Auto Inc & Khóa chính định danh tự tăng. \\ \hline
\texttt{customerId} & \texttt{VARCHAR(255)} & UK & Có & Không & Mã định danh khách hàng (UUID). \\ \hline
\texttt{name} & \texttt{VARCHAR(255)} & - & Có & Không & Họ và tên khách hàng. \\ \hline
\texttt{phoneNumber} & \texttt{VARCHAR(20)} & UK & Có & Không & Số điện thoại định danh duy nhất tại quầy. \\ \hline
\texttt{email} & \texttt{VARCHAR(255)} & - & Không & NULL & Địa chỉ email nhận hóa đơn điện tử. \\ \hline
\texttt{totalSpent} & \texttt{DOUBLE} & - & Có & 0.0 & Tổng số tiền khách hàng đã tích lũy chi tiêu. \\ \hline
\texttt{orderCount} & \texttt{INT} & - & Có & 0 & Tổng số lượng hóa đơn đã mua thành công. \\ \hline
\texttt{createdAt} & \texttt{DATETIME(6)} & - & Có & CURRENT & Thời điểm đăng ký hồ sơ khách hàng. \\ \hline
\texttt{updatedAt} & \texttt{DATETIME(6)} & - & Không & NULL & Thời điểm cập nhật hồ sơ gần nhất. \\ \hline
\end{tabularx}
\end{table}

% ------------------------------------------------------------------------------
% 11. BẢNG TBL_PROMOTIONS
% ------------------------------------------------------------------------------
\subsubsection{Bảng \texttt{tbl\_promotions} (Chương trình Khuyến mãi đa quy tắc)}
Cấu hình các chương trình ưu đãi, mã giảm giá, khung giờ vàng và mua tặng (Bảng \ref{tab:tbl_promotions}).

\begin{table}[htbp]
\centering
\caption{Từ điển dữ liệu bảng \texttt{tbl\_promotions}}
\label{tab:tbl_promotions}
\vspace{2mm}
\begin{tabularx}{\textwidth}{|p{2.5cm}|p{2.8cm}|c|c|p{1.8cm}|X|}
\hline
\textbf{Tên trường} & \textbf{Kiểu dữ liệu} & \textbf{Khóa} & \textbf{Not Null} & \textbf{Mặc định} & \textbf{Diễn giải nghiệp vụ} \\ \hline
\texttt{id} & \texttt{BIGINT} & PK & Có & Auto Inc & Khóa chính định danh tự tăng. \\ \hline
\texttt{promotionId} & \texttt{VARCHAR(255)} & UK & Có & Không & Mã định danh chương trình (UUID). \\ \hline
\texttt{name} & \texttt{VARCHAR(255)} & - & Có & Không & Tên hiển thị chương trình khuyến mãi. \\ \hline
\texttt{type} & \texttt{VARCHAR(20)} & - & Có & Không & Loại: \texttt{COUPON}, \texttt{HAPPY\_HOUR}, \texttt{BOGO}. \\ \hline
\texttt{code} & \texttt{VARCHAR(50)} & UK & Không & NULL & Mã voucher nhập tại quầy (với COUPON). \\ \hline
\texttt{discountType} & \texttt{VARCHAR(20)} & - & Không & NULL & Kiểu giảm: \texttt{PERCENTAGE} hoặc \texttt{FIXED\_AMOUNT}. \\ \hline
\texttt{discountVal} & \texttt{DECIMAL(19,2)} & - & Có & 0.00 & Giá trị chiết khấu (\% hoặc số tiền). \\ \hline
\texttt{minOrder} & \texttt{DECIMAL(19,2)} & - & Có & 0.00 & Giá trị đơn hàng tối thiểu để được áp dụng. \\ \hline
\texttt{startDate} & \texttt{DATE} & - & Không & NULL & Ngày bắt đầu có hiệu lực. \\ \hline
\texttt{endDate} & \texttt{DATE} & - & Không & NULL & Ngày kết thúc chương trình. \\ \hline
\texttt{startTime} & \texttt{TIME} & - & Không & NULL & Giờ bắt đầu áp dụng trong ngày. \\ \hline
\texttt{endTime} & \texttt{TIME} & - & Không & NULL & Giờ kết thúc áp dụng trong ngày. \\ \hline
\texttt{daysOfWeek} & \texttt{VARCHAR(255)} & - & Không & NULL & Chuỗi các ngày áp dụng trong tuần. \\ \hline
\texttt{buyVariantId} & \texttt{VARCHAR(255)} & - & Không & NULL & Mã biến thể mua điều kiện (BOGO). \\ \hline
\texttt{getVariantId} & \texttt{VARCHAR(255)} & - & Không & NULL & Mã biến thể được tặng/giảm giá (BOGO). \\ \hline
\texttt{usageLimit} & \texttt{INT} & - & Không & NULL & Giới hạn tối đa số lần sử dụng voucher. \\ \hline
\texttt{timesUsed} & \texttt{INT} & - & Có & 0 & Số lần voucher đã được áp dụng thực tế. \\ \hline
\texttt{isActive} & \texttt{TINYINT(1)} & - & Có & 1 & Trạng thái bật/tắt (1: Kích hoạt, 0: Tạm dừng). \\ \hline
\end{tabularx}
\end{table}

% ------------------------------------------------------------------------------
% 12. BẢNG TBL_ORDERS
% ------------------------------------------------------------------------------
\subsubsection{Bảng \texttt{tbl\_orders} (Hóa đơn bán hàng \& Thanh toán)}
Lưu trữ hóa đơn tổng, giá trị thanh toán, phương thức và trạng thái cổng thanh toán (Bảng \ref{tab:tbl_orders}).

\begin{table}[htbp]
\centering
\caption{Từ điển dữ liệu bảng \texttt{tbl\_orders}}
\label{tab:tbl_orders}
\vspace{2mm}
\begin{tabularx}{\textwidth}{|p{2.5cm}|p{2.8cm}|c|c|p{1.8cm}|X|}
\hline
\textbf{Tên trường} & \textbf{Kiểu dữ liệu} & \textbf{Khóa} & \textbf{Not Null} & \textbf{Mặc định} & \textbf{Diễn giải nghiệp vụ} \\ \hline
\texttt{id} & \texttt{BIGINT} & PK & Có & Auto Inc & Khóa chính, đồng thời là \texttt{orderCode} của PayOS. \\ \hline
\texttt{order\_code} & \texttt{VARCHAR(255)} & UK & Có & Không & Mã hóa đơn chuỗi hiển thị (ORD + timestamp). \\ \hline
\texttt{customerId} & \texttt{VARCHAR(255)} & - & Không & NULL & Mã khách hàng liên kết. \\ \hline
\texttt{subtotal} & \texttt{DOUBLE} & - & Có & Không & Tổng tiền hàng trước chiết khấu và thuế. \\ \hline
\texttt{discountAmt} & \texttt{DOUBLE} & - & Có & 0.0 & Số tiền được giảm từ chương trình khuyến mãi. \\ \hline
\texttt{tax} & \texttt{DOUBLE} & - & Có & 0.0 & Thuế giá trị gia tăng VAT 10\%. \\ \hline
\texttt{grandTotal} & \texttt{DOUBLE} & - & Có & Không & Tổng số tiền khách hàng thực tế phải trả. \\ \hline
\texttt{paymentMethod}& \texttt{VARCHAR(20)} & - & Có & Không & Phương thức thanh toán: \texttt{CASH} hoặc \texttt{PAYOS}. \\ \hline
\texttt{payStatus} & \texttt{VARCHAR(20)} & - & Có & PENDING & Trạng thái: PENDING, COMPLETED, CANCELLED. \\ \hline
\texttt{payQrCode} & \texttt{TEXT} & - & Không & NULL & Chuỗi dữ liệu mã VietQR động từ PayOS. \\ \hline
\texttt{createdAt} & \texttt{DATETIME(6)} & - & Có & CURRENT & Thời điểm xuất hóa đơn tại quầy. \\ \hline
\end{tabularx}
\end{table}

% ------------------------------------------------------------------------------
% 13. BẢNG TBL_ORDER_ITEMS
% ------------------------------------------------------------------------------
\subsubsection{Bảng \texttt{tbl\_order\_items} (Chi tiết mặt hàng trong hóa đơn)}
Lưu vết từng dòng sản phẩm bán ra kèm cấu hình biến thể và tùy chọn tại thời điểm xuất hóa đơn (Bảng \ref{tab:tbl_order_items}).

\begin{table}[htbp]
\centering
\caption{Từ điển dữ liệu bảng \texttt{tbl\_order\_items}}
\label{tab:tbl_order_items}
\vspace{2mm}
\begin{tabularx}{\textwidth}{|p{2.5cm}|p{2.8cm}|c|c|p{1.8cm}|X|}
\hline
\textbf{Tên trường} & \textbf{Kiểu dữ liệu} & \textbf{Khóa} & \textbf{Not Null} & \textbf{Mặc định} & \textbf{Diễn giải nghiệp vụ} \\ \hline
\texttt{id} & \texttt{BIGINT} & PK & Có & Auto Inc & Khóa chính định danh tự tăng. \\ \hline
\texttt{itemId} & \texttt{VARCHAR(255)} & - & Có & Không & Mã sản phẩm cơ sở tại thời điểm bán. \\ \hline
\texttt{variantId} & \texttt{VARCHAR(255)} & - & Không & NULL & Mã biến thể SKU vật lý xuất kho. \\ \hline
\texttt{name} & \texttt{VARCHAR(255)} & - & Có & Không & Tên sản phẩm hiển thị trên bill. \\ \hline
\texttt{basePrice} & \texttt{DOUBLE} & - & Có & Không & Giá gốc của biến thể thời điểm bán. \\ \hline
\texttt{price} & \texttt{DOUBLE} & - & Có & Không & Đơn giá cuối cùng sau khi cộng phụ thu Modifier. \\ \hline
\texttt{quantity} & \texttt{INT} & - & Có & Không & Số lượng sản phẩm bán ra. \\ \hline
\texttt{modifiers} & \texttt{TEXT} & - & Không & NULL & Danh sách JSON các tùy chọn đã chọn. \\ \hline
\texttt{order\_id} & \texttt{BIGINT} & FK & Có & Không & Khóa ngoại liên kết \texttt{tbl\_orders.id} (CASCADE). \\ \hline
\end{tabularx}
\end{table}

% ------------------------------------------------------------------------------
% 14. BẢNG TBL_ACTIVITY_LOGS
% ------------------------------------------------------------------------------
\subsubsection{Bảng \texttt{tbl\_activity\_logs} (Nhật ký kiểm toán an ninh Activity Logs)}
Lưu trữ dấu vết kiểm toán toàn diện mọi hành vi làm biến đổi dữ liệu nghiệp vụ (Bảng \ref{tab:tbl_activity_logs}).

\begin{table}[htbp]
\centering
\caption{Từ điển dữ liệu bảng \texttt{tbl\_activity\_logs}}
\label{tab:tbl_activity_logs}
\vspace{2mm}
\begin{tabularx}{\textwidth}{|p{2.5cm}|p{2.8cm}|c|c|p{1.8cm}|X|}
\hline
\textbf{Tên trường} & \textbf{Kiểu dữ liệu} & \textbf{Khóa} & \textbf{Not Null} & \textbf{Mặc định} & \textbf{Diễn giải nghiệp vụ} \\ \hline
\texttt{id} & \texttt{BIGINT} & PK & Có & Auto Inc & Khóa chính định danh tự tăng. \\ \hline
\texttt{logId} & \texttt{VARCHAR(255)} & UK & Có & Không & Mã định danh bản ghi log (UUID). \\ \hline
\texttt{userEmail} & \texttt{VARCHAR(255)} & Index & Có & Không & Email người thực hiện thao tác. \\ \hline
\texttt{action} & \texttt{VARCHAR(50)} & Index & Có & Không & Hành vi: LOGIN, CREATE, UPDATE, DELETE, CANCEL. \\ \hline
\texttt{entityType} & \texttt{VARCHAR(50)} & Index & Có & Không & Thực thể: USER, ITEM, CATEGORY, ORDER, PROMOTION... \\ \hline
\texttt{entityId} & \texttt{VARCHAR(255)} & - & Không & NULL & Mã định danh của đối tượng bị tác động. \\ \hline
\texttt{description} & \texttt{VARCHAR(1000)}& - & Không & NULL & Tóm tắt nội dung chi tiết thao tác thay đổi. \\ \hline
\texttt{createdAt} & \texttt{DATETIME(6)} & Index & Có & CURRENT & Thời điểm chính xác phát sinh thao tác. \\ \hline
\end{tabularx}
\end{table}

% ==============================================================================
\section{Tổng kết chương 2}
Chương 2 đã phân tích và mô hình hóa một cách khoa học, chuyên sâu toàn bộ bài toán nghiệp vụ của hệ thống POS bán lẻ. Thông qua hệ thống biểu đồ phân rã chức năng (BFD), mô hình Use Case tổng quát và 6 bảng đặc tả ca sử dụng chi tiết chuẩn quốc tế, các luồng tương tác giữa nhân viên thu ngân, quản trị viên và cổng thanh toán PayOS đã được định nghĩa minh bạch. Hệ thống các sơ đồ hoạt động (Activity Diagrams), tuần tự (Sequence Diagrams), sơ đồ lớp và sơ đồ triển khai phân tán đã phác họa rõ nét cách thức vận hành đồng bộ giữa các tầng kiến trúc. Đặc biệt, mô hình cơ sở dữ liệu quan hệ gồm 13 bảng chuẩn hóa và bộ Từ điển dữ liệu (Data Dictionary) chi tiết 100\% các thuộc tính, ràng buộc toàn vẹn và cơ chế kiểm toán đã tạo tiền đề kỹ thuật vững chắc để bước vào giai đoạn cài đặt, thực nghiệm và kiểm thử hệ thống trong Chương 3.

